% concatenated from arXiv_Genetic_SO_Embedding.tex
\documentclass{article}
\usepackage{amssymb}
\usepackage{amsmath}
\usepackage{amsthm}
\usepackage{color}
\usepackage[colorlinks]{hyperref}
\usepackage[a4paper]{geometry}
\usepackage{mathtools}
\usepackage{algorithm}
\usepackage{algpseudocode}
\usepackage{placeins}
\usepackage{booktabs}
\usepackage{makecell}
\usepackage{longtable}

\newtheorem{theorem}{Theorem}
\newtheorem{lemma}{Lemma}
\newtheorem{remark}{Remark}
\newcommand{\la}{\left\langle}
\newcommand{\ra}{\right\rangle}
\newcommand{\F}{\mathbb{F}}
\newcommand{\wt}{\operatorname{wt}}
\newcommand{\rank}{\operatorname{rank}}
\newcommand{\Hull}{\operatorname{Hull}}
\newcommand{\GL}{\operatorname{GL}}
\newcommand{\LCD}{\operatorname{LCD}}

\title{Constraint-Preserving Genetic Algorithms for Embedding Linear Codes into Self-Orthogonal Codes}
\author{Haeun Lim\thanks{haeunlim@sogang.ac.kr, Department of Mathematics and Institute for Mathematical and Data Sciences, Sogang University, Seoul, Korea}, Junmin An\thanks{junmin0518@sogang.ac.kr, Department of Mathematics and Institute for Mathematical and Data Sciences, Sogang University, Seoul, Korea}, Jon-Lark Kim\thanks{jlkim@sogang.ac.kr, Department of Mathematics and Institute for Mathematical and Data Sciences, Sogang University, Seoul, Korea}}
\date{}

\begin{document}
\maketitle

\begin{abstract}
In this paper, we aim to construct binary optimal self-orthogonal codes using shortest self-orthogonal embedding methods. For this purpose, we design a heuristic framework based on a genetic algorithm. We explore the search space of shortest self-orthogonal embeddings using a fitness function based on the minimum distance and the number of minimum-weight codewords. We construct \emph{constraint-preserving} crossover and mutation operations so that every chromosome yields a valid self-orthogonal embedding, while high-fitness structural features, such as favorable subsequences of orthogonal generators, are propagated across generations. We also analyze the time and storage complexity of the algorithm, and validate our design through an ablation study on guided crossover and a comparison with random search under an equal time budget. Using this method, we obtain $66$ new binary optimal self-orthogonal codes that meet the upper bound, together with $135$ further self-orthogonal codes attaining the best minimum distance found so far.\\

\textbf{Keywords} : self-orthogonal code, LCD code, embedding of codes, codes over rings
\end{abstract}


\section{Introduction}
Coding theory originated from Shannon's groundbreaking paper `A Mathematical Theory of Communication'~\cite{mat-theory-comm} and has developed into the study of error-correcting codes, which are used to detect and correct errors in noisy communication channels. One of the most fundamental problems in coding theory is to find optimal codes, namely, codes with the largest possible minimum distance for a given length and dimension. Finding optimal codes is generally a difficult problem, except for relatively small lengths and dimensions. Hence, there were many attempts to construct new optimal codes from previously known optimal codes. There are many methods that allow us to construct new codes from existing codes, including extension, puncturing, shortening, gluing~\cite{ps-1975}, Plotkin construction~\cite{p-1960}, concatenation~\cite{f-1966}, Construction X~\cite{s-1972}, matrix-product codes~\cite{bn-2001}, subfield subcodes~\cite{d-1975}, etc. By applying these methods, numerous new optimal codes have been constructed.

One of the most important classes of codes in coding theory is self-orthogonal codes. Ding~\cite{d-2009} showed that self-orthogonal codes achieve the asymptotic Gilbert-Varshamov bound. Pless and Sloane~\cite{ps-1975} constructed unimodular lattices using self-dual codes, which form a subclass of self-orthogonal codes. Calderbank et al.~\cite{crss-1998} showed that constructing quantum error-correcting codes is equivalent to finding additive self-orthogonal codes over $\mathbb{F}_4$.

In this direction, methods for constructing optimal self-orthogonal codes have been widely studied. One such well-known method is the building-up construction, which constructs a self-orthogonal $[n+2,k+1]$ code from a given self-orthogonal $[n,k]$ code~\cite{kl-2004, lk-2015, h-2021, k-2023}. Kim et al.~\cite{kll-2026} and Shi et al.~\cite{st-2026} proposed methods to expand an $[n,k]$ self-orthogonal code into an $[n,k+1]$ self-orthogonal code over finite fields $\mathbb{F}_q$ and $\mathbb{Z}_4$, respectively.

Recently, embedding method has been actively studied as a way to construct self-orthogonal codes by adding columns to a generator matrix of a given linear $[n,k]$ code. One of the central problems in embedding method is to determine the minimum number of columns required to embed a given code into a self-orthogonal code. The embedding method was first introduced in~\cite{kkl-2021}, where the authors proposed an algorithm for constructing self-orthogonal codes over $\mathbb{F}_2$ by appending the minimum number of columns to a generator matrix of a linear code of dimension $3$ or $4$. Choi et al.~\cite{kc-2022} extended the original results to dimension $6$. An et al.~\cite{akklw-2025} determined the minimum number of columns required to embed a given code into a self-orthogonal code over $\mathbb{F}_2$. Wang et al.~\cite{wl-2026} generalized these results to linear codes over arbitrary finite fields $\mathbb{F}_q$. Most recently, An et al.~\cite{akl-2026} studied the embedding of a linear code over $\mathbb{Z}_4$ into self-orthogonal codes.

Although several new optimal codes have been discovered using these embedding methods, it is almost impossible, especially for codes with large dimensions, to search through all possible embeddings in order to find optimal codes. Thus, we aim to address this problem using a heuristic algorithm. Among heuristic algorithms, we choose a genetic algorithm for the self-orthogonal embedding.

Genetic algorithms(GA) are population-based metaheuristic optimization methods inspired by the mechanisms of natural selection and evolution~\cite{h-1975}. Through gene-like operators such as crossover and mutations, GAs iteratively evolve candidate solutions toward optimality. Genetic algorithm is suitable for large-scale search problems, especially when the genetic operations can be designed to preserve the structural constraints of the problem. As a recent example of this approach in coding theory, Won et al.~\cite{w-2026} successfully developed modified genetic algorithms to systematically construct new quaternary Hermitian LCD codes. By designing custom operators—such as modified crossover and complement—that intrinsically satisfy the strict algebraic conditions of quaternary Hermitian LCD codes, their algorithms effectively navigated the highly restricted search space.

Inspired by these successful applications of evolutionary computation, we propose a tailored genetic algorithm to discover the shortest self-orthogonal embeddings with high minimum distances. Our main contribution is to reformulate the shortest self-orthogonal embedding problem into a form suitable for a genetic algorithm. We encode chromosomes using generators of the orthogonal group so that genetic operations preserve the structural constraints of self-orthogonal embedding, and then apply a genetic algorithm to explore the search space of valid embeddings. We further analyze the time and storage complexity of the algorithm, and validate our design through experiments: an ablation study confirms the benefit of the guided crossover operator, and a comparison with random search under an equal time budget demonstrates the effectiveness of the genetic algorithm. As a result, we have found $66$ new optimal self-orthogonal codes over $\mathbb{F}_2$ whose optimal minimum distances were previously not known, together with $135$ further self-orthogonal codes attaining the best minimum distance obtained so far.

The remainder of this paper is organized as follows. Section 2 reviews the basic notions of coding theory and genetic algorithms. Section 3 formulates the shortest self-orthogonal embedding problem, presents our constraint-preserving genetic algorithm, and analyzes its complexity. Section 4 reports the numerical results, together with the running time, an ablation study, and a comparison with random search. Finally, Section 5 concludes the paper.

\section{Preliminaries}
\subsection{Introduction to Coding Theory}
Let $\mathbb{F}_2$ be the binary field with two elements. A \textit{code} $\mathcal{C}$ over $\mathbb{F}_2$ is a subset of $\mathbb{F}_2^n$ for $n \geq 1$. In this case, $n$ is the length of $\mathcal{C}$. Elements of $\mathcal{C}$ are called \textit{codewords}. A \textit{binary linear code} is a $k$-dimensional subspace of $\mathbb{F}_2^n$ and denoted by an $[n,k]$ code. We only consider binary linear codes in this paper. A \textit{generator matrix} $G$ for an $[n,k]$ code $\mathcal{C}$ is a binary $k \times n$ matrix whose rows form a basis for $\mathcal{C}$. For an $[n,k]$ code  $\mathcal{C}$, there is a generator matrix $G$ such that
\[
G = [~I_k~~~A~]
\]
where $I_k$ is the $k \times k$ identity matrix. Such a generator matrix is called a \textit{standard generator matrix} for $\mathcal{C}$ and is denoted by $G(\mathcal{C})$. Two codes $\mathcal{C}_1$ and $\mathcal{C}_2$ over $\mathbb{F}_2$ are called \textit{equivalent} if there exists a permutation of columns $\sigma$ where $\sigma \mathcal{C}_1 = \mathcal{C}_2$. The \textit{permutation automorphism group} of $\mathcal{C}$, denoted by $\mathrm{Aut}(\mathcal{C})$, is the group of all coordinate permutations that fix $\mathcal{C}$, that is, $\mathrm{Aut}(\mathcal{C}) = \{\sigma \in S_n ~|~ \sigma\mathcal{C} = \mathcal{C}\}$, where $S_n$ is the symmetric group on the $n$ coordinate positions. Its order is denoted by $|\mathrm{Aut}(\mathcal{C})|$. Since equivalent codes have isomorphic permutation automorphism groups, the order $|\mathrm{Aut}(\mathcal{C})|$ is an invariant of the equivalence class of $\mathcal{C}$. Hence two codes with the same parameters $[n,k,d]$ but distinct automorphism group orders are inequivalent, which we use to distinguish the codes obtained in Section 4.
The \textit{dual} $\mathcal{C}^{\perp}$ of a code $\mathcal{C}$ is defined as the set $\{\mathbf{x} \in \mathbb{F}_2^n~|~\mathbf{x} \cdot \mathbf{c} = 0  \textrm{ for all } \mathbf{c} \in \mathcal{C}$ \} where $\cdot$ is the standard dot product. For an $[n, k]$ code $\mathcal{C}$ over $\mathbb{F}_2$, $\dim \mathcal{C}^{\perp} = n-k$. If $\mathcal{C} \subseteq \mathcal{C}^{\perp}$, then $\mathcal{C}$ is called a \textit{self-orthogonal code}. Moreover, if $\mathcal{C} = \mathcal{C}^\perp$, then it is called a \textit{self-dual code}. Given any code $\mathcal{C}$, the dual of $\mathcal{C}$ is a linear code. Consequently, a self-dual code is a linear code. The \textit{hull} of $\mathcal{C}$, denoted by $\mathrm{Hull}(\mathcal{C})$ is defined as \[
\mathrm{Hull}(\mathcal{C}) = \mathcal{C} \cap \mathcal{C}^\perp.
\]
In particular, $\mathcal{C}$ is self-orthogonal if and only if $GG^{T}=\mathcal{O}$.
The \textit{Hamming weight} of a vector $\mathbf{x} \in \mathbb{F}_2^n$, denoted by $\wt(\mathbf{x})$, is the number of 1's in $\mathbf{x}$. For two vectors $\mathbf{x}, \mathbf{y}$ in $\mathbb{F}_2^n$, the \textit{Hamming distance} between $\mathbf{x}$ and $\mathbf{y}$ is defined as
\[
d(\mathbf{x}, \mathbf{y})=\wt(\mathbf{x} - \mathbf{y}).
\]
Moreover, the weight of the sum of two vectors $\mathbf{x}, \mathbf{y}$ is defined as \[
\wt(\mathbf{x}+ \mathbf{y})=\wt(\mathbf{x})+\wt(\mathbf{y})-2(\mathbf{x} * \mathbf{y})
\] where $\mathbf{x}*\mathbf{y}$ denotes the number of coordinates in which both vectors have entry 1.
The \textit{minimum (Hamming) distance} of a code $\mathcal{C}$, denoted by $d(\mathcal{C})$, is the minimum of the distances between any two distinct codewords in $\mathcal{C}$. An $[n,k]$ code $\mathcal{C}$ with minimum distance $d$ is denoted by an $[n,k,d]$ code. For linear codes, minimum distance equals the \textit{minimum (Hamming) weight} among its nonzero codewords in $\mathcal{C}$. For a code $\mathcal{C}$ over $\mathbb{F}_2$, if every codeword of $\mathcal{C}$ has even weight, $\mathcal{C}$ is called an \textit{even code}. Otherwise, if it contains an odd-weight codeword, it is called an \textit{odd-like code}.
Given $n$ and $k$, we denote by $d(n, k)$ the largest possible minimum distance over all linear $[n, k]$ codes. A linear $[n,k]$ code $\mathcal{C}$ with minimum distance $d(\mathcal{C})$ is called \textit{optimal} if
\[
d(\mathcal{C})=d(n,k).
\]
Similarly, $d_{SO}(n, k)$ is the largest possible minimum distance over all self-orthogonal $[n, k]$ codes. A self-orthogonal $[n,k]$ code $\mathcal{C}$ is called an \textit{optimal self-orthogonal code} if
\[
d(\mathcal{C}) = d_{SO}(n, k).
\]
For an $[n,k]$ code $\mathcal{C}$, let $S$ be a subset of the set of coordinate positions of $\mathcal{C}$. The code obtained by removing the columns indexed by $S$ from a generator matrix $G(\mathcal{C})$ is called a \textit{punctured code} of $\mathcal{C}$. The punctured code has length $n-m$, where $m = |S|$.
Given an $[n, k, d]$ code $\mathcal{C}$ and $\mathbf{c}$, a vector of weight $d$ in $\mathcal{C}$, define the residual code of $\mathcal{C}$ with respect to $\mathbf{c}$, denoted by $\mathrm{Res}(\mathcal{C}, \mathbf{c})$, as the code of length $n-d$, obtained by puncturing all coordinate positions where $\mathbf{c}$ has ones.
\begin{theorem}[{\cite{Introduction}}]
If $\mathcal{C}$ is an $[n, k, d]$ code over $\mathbb{F}_2$ and $\mathbf{c}$ is a codeword in $\mathcal{C}$ of weight $d$, then $\mathrm{Res}(\mathcal{C}, \mathbf{c})$ is an $[n-d, k-1, d']$ code where $\displaystyle d' \geq \left\lceil\frac{d}{2}\right\rceil$.
\end{theorem}
\subsection{Genetic Algorithm}
A genetic algorithm is a population-based heuristic search algorithm inspired by the mechanisms of natural evolution and adaptation. It is used to find good approximate solutions to optimization problems, especially when the search space is too large for exhaustive search. In a genetic algorithm, candidate solutions are represented as artificial \textit{chromosomes}, and each chromosome is evaluated by a \textit{fitness function}, which measures how promising the chromosome is as a solution. For a comprehensive treatment of genetic algorithms, we refer the reader to~\cite{goldberg-1989}.
The basic idea of a genetic algorithm is that high-quality ``parent" candidate solutions from different regions in the space may produce high-quality ``offspring" candidate solutions. Based on this idea, chromosomes with higher fitness values are more likely to be selected as parents. The selected parents then generate new chromosomes through \textit{crossover}, which combines parts of two parent chromosomes. In addition, \textit{mutation} is applied to introduce small random changes to the offspring chromosomes and to maintain diversity in the population.
We now describe the main steps of a genetic algorithm in more detail.

\textbf{Chromosomes and genes.}
A chromosome represents a candidate solution to the given problem. It is often encoded as a bit string, although the encoding may vary depending on the problem. The basic components of a chromosome are called \textit{genes}. In a bit-string representation, genes may be single bits or short blocks of adjacent bits encoding particular elements of the candidate solution.

\textbf{Selection.}
Selection is the process of choosing chromosomes from the current population to serve as parents for the next generation. We also use \textit{elitism}, which preserves the best chromosomes from the previous generation. After preserving the elite chromosomes, the remaining positions in the next generation are filled through selection. For the parent selection mechanism, we employ \textit{tournament selection}. In each selection event, $t$ candidate chromosomes are randomly drawn from the population, and the individuals with high fitness function values among them are selected. This procedure is iteratively executed until the required number of parents is collected.

\textbf{Crossover.}
Crossover combines two parent chromosomes to produce offspring. In one-point crossover, a crossover point is chosen, and the subsequences before and after this point are exchanged between the two parents. For example, if the two parent strings are
\[
10001000
\qquad\text{and}\qquad
11111111,
\]
then crossing over after the third position produces
\[
10011111
\qquad\text{and}\qquad
11101000.
\]
This operation roughly mimics biological recombination. If each pair of parents produces two offspring and the population is replaced by the same number of chromosomes, then the population size remains unchanged.

\textbf{Mutation.}
Mutation introduces random changes into a chromosome. In the bit-string representation, mutation is often implemented by flipping randomly selected bits. More generally, mutation operations can be designed according to the structure of the problem.

The sequence of selection, crossover, and mutation is iterated for a predetermined number of \textit{generations}, or until a prescribed stopping criterion is satisfied.
In many applications of genetic algorithms, chromosomes can be generated and modified rather freely. However, in some problems, not every formal chromosome represents a feasible solution. In such cases, the chromosome must satisfy certain structural constraints, and the genetic operations must be designed so that these constraints are preserved after crossover and mutation.

\section{Problem Formulation and a Constraint-Preserving Genetic Algorithm}
In this section, we introduce the shortest self-orthogonal embedding problem and present a customized genetic algorithm suitable for our problem. Specifically, we construct an algorithmic framework equipped with constraint-preserving operators to efficiently navigate this restricted search space.
\subsection{Shortest self-orthogonal embeddings and the feasibility condition}
Given an $[n, k]$ code $\mathcal{C}$ over $\mathbb{F}_2$, an $[n', k]$ self-orthogonal code $\mathcal{C}'$ is called a \textit{self-orthogonal embedding} of $\mathcal{C}$ over $\mathbb{F}_2$ if $\mathcal{C}$ is obtained by puncturing $\mathcal{C}'$. Let $m = n'-n$. Then there exists a set $\mathcal{I}=\{i_1,\ldots,i_m\}$ of coordinate positions of $\mathcal{C}'$ such that puncturing $\mathcal{C}'$ at the coordinates in $\mathcal{I}$ gives $\mathcal{C}$. Up to equivalence, we may assume that the columns indexed by $\mathcal{I}$ are placed in the last $m$ coordinate positions. Thus, for a self-orthogonal embedding $\mathcal{C}'$ of $\mathcal{C}$, we consider a generator matrix of $\mathcal{C}'$ of the form \[
G(\mathcal{C}') = [G(\mathcal{C})~|~A].
\]
For a code $\mathcal{C}$ over $\mathbb{F}_2$, consider the code generated by $G' = [G(\mathcal{C})~|~G(\mathcal{C})]$. Since
\[
G'G'^T = G(\mathcal{C})G(\mathcal{C})^T + G(\mathcal{C})G(\mathcal{C})^T = \mathcal{O},
\]
the code generated by $G'$ is self-orthogonal. It follows that any linear code over $\mathbb{F}_2$ has a self-orthogonal embedding.
In coding theoretic perspective, codes with smaller length are often more efficient. Hence, it is natural to ask for the minimum number of columns that must be added to obtain a self-orthogonal embedding. A \textit{shortest self-orthogonal embedding} of $\mathcal{C}$ is a self-orthogonal embedding of $\mathcal{C}$ whose length is the smallest among all self-orthogonal embeddings of $\mathcal{C}$. Using the notion of the hull of a code $\mathcal{C}$, An et al. determined the minimum number of columns needed to embed $\mathcal{C}$ into a self-orthogonal code as follows.
\begin{theorem}\cite{akklw-2025}\label{akklw-2025}
    Let $\mathcal{C}$ be a binary $[n,k]$ code and let $\ell=\dim \Hull(\mathcal{C})$. Then the length $n'$ of a shortest self-orthogonal embedding of $\mathcal{C}$ is
    \[
    n' = \begin{cases}
        n+k-\ell+1,   &\mathrm{if~}\mathcal{C}\textup{~is~even,}\\
        n+k-\ell,     &\mathrm{if~}\mathcal{C}\textup{~is~odd-like.}
    \end{cases}
    \]
\end{theorem}
Let $\mathcal{C}$ be an $[n,k]$ code over $\mathbb{F}_2$ with $\dim\Hull(\mathcal{C}) = \ell$. Let $H$ be a generator matrix of $\Hull(\mathcal{C})$. Since $\Hull(\mathcal{C})$ is a subcode of $\mathcal{C}$, we can extend $H$ to a generator matrix $G(\mathcal{C})$ of $\mathcal{C}$ as follows:
\[
G(\mathcal{C}) = \begin{bmatrix}
    H \\ A
\end{bmatrix}
\]
for some $(k-\ell)\times n$ matrix $A$.
% \begin{comment} Moreover,
% \begin{align*}
%     k-\ell &= \rank(G(\mathcal{C})G(\mathcal{C})^T)
%     =\rank\begin{bmatrix}
%         HH^T & HA^T \\
%         AH^T & AA^T
%     \end{bmatrix}\\
%     &=\rank\begin{bmatrix}
%     \mathcal{O} & \mathcal{O}\\
%     \mathcal{O} & AA^T
%     \end{bmatrix}
%     =\rank(AA^T),
% \end{align*}
% where $\mathcal{O}$ is the zero matrix.
% Hence, by Equation~\ref{hulldimformula}, the code generated by $A$ is an LCD code.\end{comment}
\begin{theorem}\label{zero-embedding}\cite{akklw-2025}
Let $\mathcal{C}$ be an $[n,k]$ code over $\mathbb{F}_2$ with $\ell=\dim\Hull(\mathcal{C})$. Let \[
G(\mathcal{C}) = \begin{bmatrix}
    H \\ A
\end{bmatrix}
\] where $H$ is a generator matrix of $\Hull(\mathcal{C})$. Then, a shortest self-orthogonal embedding $\mathcal{C}'$ of $\mathcal{C}$ has a generator matrix
\[
G(\mathcal{C}')=\begin{bmatrix}
    H & \mathcal{O}\\
    A & B
    \end{bmatrix}
\] for some matrix $B$.
\end{theorem}
By Theorem~\ref{zero-embedding}, every shortest self-orthogonal embedding of $\mathcal{C}$ can be obtained, up to equivalence, by varying only the block $B$ in $G(\mathcal{C}')$. Note that
\begin{align*}
    G(\mathcal{C}')G(\mathcal{C}')^T &= \begin{bmatrix}
    H & \mathcal{O}\\
    A & B
    \end{bmatrix}\begin{bmatrix}
    H^T & A^T\\
    \mathcal{O} & B^T
    \end{bmatrix}\\&=\begin{bmatrix}
    HH^T & HA^T\\
    AH^T & AA^T+BB^T
    \end{bmatrix}\\&=\begin{bmatrix}
    \mathcal{O} & \mathcal{O}\\
    \mathcal{O} & AA^T+BB^T
    \end{bmatrix}.
\end{align*}
Since $\mathcal{C}'$ is self-orthogonal, we have
\[
G(\mathcal{C}')G(\mathcal{C}')^T=\mathcal{O}.
\]
By the above computation, this condition reduces to
\[
AA^T+BB^T=\mathcal{O}.
\]
Equivalently, over $\mathbb{F}_2$, the block $B$ must satisfy
\[
BB^T=AA^T.
\]
% O(n, 2) matrix를 B에 곱해서 모든 B generation 가능
% orthogonal matrix의 generator 설계
% genetic algorithm 적용
\subsection{A constraint-preserving genetic algorithm}\label{so-embedding-algorithm}
An et al.~\cite{akklw-2025} showed that once a generator matrix of a shortest self-orthogonal embedding is obtained, all other shortest self-orthogonal embeddings can be constructed using the following theorem.
\begin{theorem}\cite{akklw-2025}\label{orthogonal-matrix}
    Let $\mathcal{C}$ be an $[n,k]$ code over $\mathbb{F}_2$ with a generator matrix \[
    G(\mathcal{C}) = \begin{bmatrix}
        H \\ A
    \end{bmatrix}
    \] where $H$ is a generator matrix of $\Hull(\mathcal{C})$.
    Let $\mathcal{C}'$ be an $[n', k]$ shortest self-orthogonal embedding of $\mathcal{C}$ with a generator matrix\[
    G(\mathcal{C}')=\begin{bmatrix}
        H & \mathcal{O}\\
        A & B
        \end{bmatrix}.
    \]
    Let $m=n'-n$. A code $\tilde{\mathcal{C}}$ is a shortest self-orthogonal embedding of $\mathcal{C}$ if and only if it has a generator matrix \[
    G(\tilde{\mathcal{C}}) = \begin{bmatrix}
    H & \mathcal{O}\\
    A & BD
    \end{bmatrix}
\] for $D \in O(m,2)$ where $O(m,2)$ is the orthogonal group of order $m$ over $\mathbb{F}_2$ given as follows:
\[
O(m,2) = \{M \in GL_m(\mathbb{F}_2)~|~MM^T=I \}.
\]
\end{theorem}
For $n \geq 1$, let $\mathbf{u}$ be a vector in $\mathbb{F}_2^m$. We define the \textit{transvection} $T_\mathbf{u}:\mathbb{F}_2^m \to \mathbb{F}_2^m$ as \[
T_\mathbf{u}(\mathbf{x})=\mathbf{x}+(\mathbf{x} \cdot \mathbf{u})\mathbf{u}=\mathbf{x}(I_m+\mathbf{u}^T\mathbf{u})
\] for every $\mathbf{x} \in \mathbb{F}_2^m$. From now on, we identify $T_\mathbf{u}$ with $I_m+\mathbf{u}^T\mathbf{u}$.
\begin{remark}\label{orthogonal-group}
    It is known from \cite{Parametrization} that the orthogonal group $O(m,2)$ of order $m$ over $\mathbb{F}_2$ is given as follows:
    \begin{enumerate}
    \item if $1 \leq m \leq 3$, then $O(m,2) = \mathcal{P}_m$ where $\mathcal{P}_m$ is a permutation group of order $m$,
    \item if $m \geq 4$, then $O(m,2)=\langle\mathcal{P}_m, T_{\mathbf{u}}\rangle$
    where $\mathbf{u}$ is a vector of Hamming weight 4 over $\mathbb{F}_2$.
    \end{enumerate}
\end{remark}
\begin{theorem}\label{orthogonal-group-generator}
    For $m \geq 5$, $O(m,2)$ over $\mathbb{F}_2$ is generated by \[
    \mathcal{T} = \{T_\mathbf{u}~|~\wt(\mathbf{u})=4~for~\mathbf{u}\in\mathbb{F}_2^m\}.
    \]
\end{theorem}
\begin{proof}
    For any $\mathbf{u}_0 \in \mathbb{F}_2^m$, by Remark~\ref{orthogonal-group}, $O(m,2)=\langle\mathcal{P}_m, T_{\mathbf{u}_0}\rangle$. Thus, it is straightforward that $T_{\mathbf{u}_0} \in O(m,2)$.
    Next, we show that $\mathcal{T}$ generates $O(m,2)$, i.e., $\mathcal{P}_m \subseteq \langle\mathcal{T}\rangle$. Since $\mathcal{P}_m$ is generated by transpositions, it is sufficient to prove that any transposition $(ab)$ is generated by transvections in $\mathcal{T}$. We will show that any transposition can be written as a product of three transvections. Choose three distinct indices $c, d, e \notin \{a, b\}$, and set\[
    p=\mathbf{e}_a + \mathbf{e}_b + \mathbf{e}_d + \mathbf{e}_e,~~~~q=\mathbf{e}_b + \mathbf{e}_c + \mathbf{e}_d + \mathbf{e}_e.
    \] where $\mathbf{e}_i$ denotes the $i$-th standard basis vector of $\mathbb{F}_2^m$, that is, the vector whose $i$-th coordinate is $1$ and all other coordinates are $0$.
    Then $\wt(p)=\wt(q)=4$ and $p \cdot q = 1$. For $\mathbf{x} \in \mathbb{F}_2^m$, we have
\begin{align*}
    T_p T_q T_p(\mathbf{x}) &= T_p T_q(\mathbf{x} + (\mathbf{x} \cdot p)p)\\
    &= T_p (\mathbf{x} + (\mathbf{x} \cdot p)p + (\mathbf{x} \cdot q + \mathbf{x} \cdot p)q)\\
    &= \mathbf{x} + (\mathbf{x} \cdot p)p + (\mathbf{x} \cdot q + \mathbf{x} \cdot p)q + (\mathbf{x} \cdot q)p\\
    &= \mathbf{x} + (\mathbf{x} \cdot (p+q))(p+q)\\
    &= T_{p+q}(\mathbf{x})\\
    &= T_{\mathbf{e}_a+\mathbf{e}_b}(\mathbf{x}).
\end{align*}
For $\mathbf{x} \in \mathbb{F}_2^m$,
\begin{align*}
    T_{\mathbf{e}_a+\mathbf{e}_b}(\mathbf{x})&=\mathbf{x} + (\mathbf{x} \cdot (\mathbf{e}_a+\mathbf{e}_b))(\mathbf{e}_a+\mathbf{e}_b)\\
    &=\mathbf{x} + (x_a+x_b)(\mathbf{e}_a+\mathbf{e}_b)\\
    &=\mathbf{x} + (x_a+x_b)\mathbf{e}_a + (x_a+x_b)\mathbf{e}_b\\
\end{align*}
Hence, we conclude that $T_{\mathbf{e}_a+\mathbf{e}_b}$ interchanges the $a$-th and $b$-th coordinates of $\mathbf{x}$ and fixes all other coordinates.
Thus, $T_p T_q T_p = T_{\mathbf{e}_a+\mathbf{e}_b} = (ab)$.
\end{proof}
\begin{remark}[\cite{m-1969}]\label{orthogonal-group-order}
The exact order of the orthogonal group $O(m,2)$ is known. Let $t=\lfloor m/2\rfloor$. Then
\[
|O(m,2)|=
\begin{cases}
\displaystyle 2^{t}\prod_{i=1}^{t-1}\bigl(2^{2t}-2^{2i}\bigr), & \text{if } m=2t,\\[3mm]
\displaystyle 2^{t}\prod_{i=0}^{t-1}\bigl(2^{2t}-2^{2i}\bigr), & \text{if } m=2t+1.
\end{cases}
\]
In particular, $|O(2t+1,2)|=(2^{2t}-1)\,|O(2t,2)|$. Table~\ref{tab:orders-orthogonal-groups} lists these orders for small values of $m$.
\end{remark}
\begin{table}[h]
\centering
\begin{tabular}{c|c}
\hline
$m$ & $|O(m, 2)|$ \\
\hline
$4$ & $2\cdot 4!$ \\
$5$ & $6\cdot 5!$ \\
$6$ & $2^5\cdot 6!$ \\
$7$ & $9\cdot 2^5\cdot 7!$ \\
$8$ & $9\cdot 2^9\cdot 8!$ \\
$9$ & $255\cdot 2^9\cdot 9!$ \\
\hline
\end{tabular}
\caption{Orders of $O(m, 2)$ for small values of $m$.}
\label{tab:orders-orthogonal-groups}
\end{table}
Let $\mathcal{C}$ be an $[n,k]$ code over $\mathbb{F}_2$. Our goal is to find a self-orthogonal embedding of $\mathcal{C}$ with the largest possible minimum distance. Let $\mathcal{C}'$ be an $[n+m,k]$ self-orthogonal embedding of $\mathcal{C}$ over $\mathbb{F}_2$. Suppose that
\[
G(\mathcal{C})=
\begin{bmatrix}
H\\
A
\end{bmatrix}
\qquad \text{and} \qquad
G(\mathcal{C}')=
\begin{bmatrix}
H & \mathcal{O}\\
A & B_0
\end{bmatrix},
\]
where $B_0$ is a matrix satisfying
\[
AA^T=B_0B_0^T.
\]
In order to find all shortest self-orthogonal embeddings of $\mathcal{C}$, we need to compute the minimum distance of the code generated by
\[
\begin{bmatrix}
H & \mathcal{O}\\
A & B_0D
\end{bmatrix}
\]
for every matrix $D\in O(m,2)$. However, as shown in Table~\ref{tab:orders-orthogonal-groups} and Remark~\ref{orthogonal-group-order}, an exhaustive search over all matrices in the orthogonal group is computationally infeasible.
Here, we apply a genetic algorithm to search for self-orthogonal embeddings, based on Theorem~\ref{orthogonal-group-generator}. Assume that the number of appended columns is $m\geq 5$. Let
\[
U=\{\mathbf{u}_1,\mathbf{u}_2,\ldots,\mathbf{u}_r\}
\]
be the set of all weight-$4$ vectors in $\mathbb{F}_2^m$.
We define a chromosome to be a finite sequence of vectors from $U$. More precisely, for a positive integer $s$, a chromosome is a sequence
\[
\chi=(\mathbf{u}_{i_1},\mathbf{u}_{i_2},\ldots,\mathbf{u}_{i_s}),
\]
where $\mathbf{u}_{i_j}\in U$ for each $j=1,\ldots,s$. We define the length of $\chi$ to be the number of transvections represented by the sequence, namely
\[
|\chi|=s.
\]
For each chromosome $\chi$, we associate the matrix
\[
D_{\chi}
=
T_{\mathbf{u}_{i_1}}
T_{\mathbf{u}_{i_2}}
\cdots
T_{\mathbf{u}_{i_s}}.
\]
By Theorem~\ref{orthogonal-group-generator}, $D_{\chi}\in O(m,2)$.
Before describing the genetic operations, we define the fitness function used to evaluate a chromosome. Let $\chi$ be a chromosome and let $D_{\chi}$ be the corresponding orthogonal matrix. We denote by $\mathcal{C}_{\chi}$ the code generated by
\[
G_{\chi}=
\begin{bmatrix}
H & \mathcal{O}\\
A & B_0D_{\chi}
\end{bmatrix}.
\]
The primary objective is to maximize the minimum distance of $\mathcal{C}_{\chi}$. However, many chromosomes may yield codes with the same minimum distance. To distinguish such chromosomes, we use a \textit{lexicographic fitness function}. Let
\[
d_{\chi}=d(\mathcal{C}_{\chi})
\]
be the minimum distance of $\mathcal{C}_{\chi}$, and let
\[
A_i(\mathcal{C}_{\chi})
=
\left|\left\{
\mathbf{c}\in \mathcal{C}_{\chi}:\operatorname{wt}(\mathbf{c})=i\right\}\right|
\]
be the number of codewords of weight $i$ in $\mathcal{C}_{\chi}$. We define the \textit{fitness} of $\chi$ by
\[
F(\chi)
=
\left(
d_{\chi},
-A_{d_{\chi}}(\mathcal{C}_{\chi}),
-|\chi|
\right),
\]
where $|\chi|$ denotes the length of the chromosome. Fitness values are compared lexicographically. Thus, a chromosome with larger minimum distance is preferred. If two chromosomes have the same minimum distance, then the one with fewer minimum-weight codewords is preferred. The last component is used only as a tie-breaker to discourage unnecessarily long representations.
Even when two chromosomes yield codes with the same minimum distance, the lexicographic fitness function allows us to select a more promising chromosome among candidates.
In this algorithm, there are two types of crossovers, random crossover and guided crossover.
Let
\[ \chi_1=(\mathbf{u}_{i_1}, \mathbf{u}_{i_2}, \ldots, \mathbf{u}_{i_s}) \qquad \text{and} \qquad \chi_2=(\mathbf{u}_{j_1}, \mathbf{u}_{j_2}, \ldots, \mathbf{u}_{j_{s'}}).\]
First, we introduce the random crossover operation. Crossover points $a$ and $b$ are selected randomly within $1 \leq a \leq s$ and $1 \leq b \leq s'$. Then the output chromosome is \[
\chi_3 = (\mathbf{u}_{i_1}, \ldots, \mathbf{u}_{i_a},\mathbf{u}_{j_{b+1}}, \ldots, \mathbf{u}_{j_{s'}}),
\]
and
\[
\chi_4 = (\mathbf{u}_{j_1}, \ldots, \mathbf{u}_{j_b},\mathbf{u}_{i_{a+1}}, \ldots, \mathbf{u}_{i_{s}}).
\]
% block crossover
Next, we introduce the guided crossover operation. We regard $\chi_1$ as the base chromosome and $\chi_2$ as the secondary chromosome. Let $\chi_1=(\mathbf{u}_{i_1},\mathbf{u}_{i_2},\ldots,\mathbf{u}_{i_s})$ be the base chromosome. For each $1\leq r\leq s$, define
\[
{\chi_1}_r=(\mathbf{u}_{i_1},\ldots,\mathbf{u}_{i_r}).
\]
For a positive integer $p\leq s$, we select the $p$ \textit{prefix chromosomes} among $\chi_{1_1},\chi_{1_2},\ldots,\chi_{1_s}$ with the largest fitness values. Similarly, let $\chi_2=(\mathbf{u}_{j_1},\mathbf{u}_{j_2},\ldots,\mathbf{u}_{j_{s'}})$ be the secondary chromosome. For each $1\leq r'\leq s'$, define
\[
{\chi_2}^{r'}=(\mathbf{u}_{j_{r'}},\ldots,\mathbf{u}_{j_{s'}}).
\]
For a positive integer $q\leq s'$, we select the $q$ \textit{suffix chromosomes} among ${\chi_2}^1,{\chi_2}^2,\ldots,{\chi_2}^{s'}$ with the largest fitness values. We then concatenate each selected chromosome from the base chromosome with each selected chromosome from the secondary chromosome. In this way, we obtain $pq$ candidate offspring chromosomes. Among these candidates, we choose two chromosomes $\chi_3$ and $\chi_4$ with the largest fitness values and use them as offspring.
To prevent excessive chromosome length, we constrain the maximum length of the resulting chromosomes for both crossover operators. The maximum allowable length of a chromosome is set as a hyperparameter, denoted by $L_{\max}$. In our experiment, we set $L_{\max}=4m$. If $|\chi_3|>L_{\max}$ or $|\chi_4|>L_{\max}$, then the crossover points are chosen again.
There are several types of mutation operations. Let
\[
\chi=(\mathbf{u}_{i_1},\mathbf{u}_{i_2},\ldots,\mathbf{u}_{i_s})
\]
be a chromosome. We first introduce four mutation operations.
\begin{enumerate}
    \item \textbf{Replace.} Replace a randomly selected entry $\mathbf{u}_{i_\ell}$ of $\chi$ with a new vector $\mathbf{u}_{i_t}\in U$.

    \item \textbf{Insert.} Insert a new vector $\mathbf{u}_{i_t}\in U$ at a randomly selected position in $\chi$.

    \item \textbf{Delete.} Delete a randomly selected entry $\mathbf{u}_{i_\ell}$ from $\chi$.

    \item \textbf{Swap.} Interchange two randomly selected entries $\mathbf{u}_{i_{\ell_1}}$ and $\mathbf{u}_{i_{\ell_2}}$ of $\chi$. Since the orthogonal group $O(m,2)$ is generally non-abelian, changing the order of transvections may produce a different corresponding orthogonal matrix.
\end{enumerate}
After crossover and mutation, each offspring is refined by a bounded local search: for up to $D$ trials, a random weight-four transvection $\mathbf{u}$ is appended to the chromosome and kept only if it strictly improves the fitness, stopping early once the length bound $L_{\max}=4m$ is reached. This refinement is applied to every offspring but not to the elite chromosomes. In our experiments we set $D=8$.
\subsection{Complexity}\label{sec:complexity}
We analyze the cost of evaluating the fitness of a candidate embedding, which is the dominant operation of the algorithm. Throughout this subsection we write $B$ for the appended block $B_0D_{\chi}$ of a candidate, since from a computational standpoint only the resulting block matters.

Let $\mathcal{C}$ be the base $[n,k,d]$ code to be embedded, and let $\ell=\dim\Hull(\mathcal{C})$ and $r=k-\ell$. Write a generator matrix of $\mathcal{C}$ as
\[
G(\mathcal{C})=\begin{bmatrix} H \\ A \end{bmatrix},
\]
where the rows of $H$ form a basis of $\Hull(\mathcal{C})$ and the rows $a_1,\ldots,a_r$ of $A$ complete it to a basis of $\mathcal{C}$. Then every codeword $c\in\mathcal{C}$ is written uniquely as
\[
c=h+a_x,\qquad h\in\Hull(\mathcal{C}),\quad a_x:=x_1a_1+\cdots+x_ra_r,
\]
for some $x=(x_1,\ldots,x_r)\in\mathbb{F}_2^r$. For a fixed $x$, the codewords $a_x+h$ with $h\in\Hull(\mathcal{C})$ form the coset $a_x+\Hull(\mathcal{C})$.

For a candidate embedding with appended block $B$, the generator matrix has the form $\begin{bmatrix} H & \mathcal{O} \\ A & B \end{bmatrix}$, so the appended coordinates of the codeword $h+a_x$ are $b_x:=x_1b_1+\cdots+x_rb_r$, where $b_i$ is the $i$-th row of $B$. Crucially, $b_x$ depends only on $x$ and not on $h$: within a coset the appended part is constant. Since the original and appended coordinates are disjoint, the weight splits as
\[
\operatorname{wt}(h+a_x,\,b_x)=\operatorname{wt}(h+a_x)+\operatorname{wt}(b_x),
\]
and the minimum weight over the coset is therefore
\[
\min_{h\in\Hull(\mathcal{C})}\operatorname{wt}(h+a_x,\,b_x)
=\Big(\min_{h\in\Hull(\mathcal{C})}\operatorname{wt}(h+a_x)\Big)+\operatorname{wt}(b_x).
\]

The first term does not depend on the candidate $B$. Hence, once per base code, we precompute the values
\[
\mu_x=\min_{h\in\Hull(\mathcal{C})}\operatorname{wt}(a_x+h)
\qquad (x\in\mathbb{F}_2^r).
\]
This is implemented by enumerating all $2^{\ell}$ hull words for each of the $2^{r}$ cosets, and therefore costs $O(2^{\ell}2^{r})=O(2^{k})$ time. Given any candidate $B$, its minimum distance is then obtained as
\[
\min_{x\in\mathbb{F}_2^r}\big(\mu_x+\operatorname{wt}(b_x)\big),
\]
which scans only the $2^{r}$ cosets rather than all $2^{k}$ codewords, i.e.\ in $O(2^{r})$ time.

Consequently, let $P$, $G$, $D$, and $L$ denote the population size, number of generations, local-search depth, and maximum chromosome length, respectively. In each generation, the algorithm evaluates the fitness of $O(P)$ candidates, and each candidate undergoes up to $D$ local-search steps that each trigger a fitness evaluation; over $G$ generations this amounts to
\[
E=O(GPD)
\]
fitness evaluations in total. Each evaluation costs $O(2^{r})$ by the coset scan above, and the one-time precomputation costs $O(2^{k})$, so the total running time is
\[
O\!\big(2^{k}+E\,2^{r}\big)=O\!\big(2^{k}+GPD\,2^{r}\big).
\]
For the storage, the precomputed table $\{\mu_x\}$ requires $O(2^{r})$ space, and maintaining a population of $P$ chromosomes, each a gene sequence of length at most $L$, requires $O(PL)$ space; hence the total storage is $O\!\big(2^{r}+PL\big)$. In contrast, an exhaustive search over the orthogonal group would evaluate all $|O(m,2)|=2^{\Theta(m^{2})}$ matrices. Since $m=k-\ell$ or $k-\ell+1$, we have $m\le k$, so the exhaustive cost $2^{\Theta(m^{2})}$ grows with the \emph{square} of the exponent, whereas the cost of our algorithm is governed by $2^{k}$, whose exponent is only \emph{linear} in $k$. Although both are exponential in the code dimension, reducing the exponent from order $m^{2}$ to order $k$ is a substantial saving: for instance, when $\ell$ is small so that $m\approx k$, the exhaustive search scales as $2^{\Theta(k^{2})}$ while our algorithm scales as $2^{\Theta(k)}$. The genetic algorithm thus replaces the search over a group of size $2^{\Theta{(m^2)}}$ by a fixed number of evaluations, each costing only $O(2^{r})$.
% guided mutation
% random kick
\section{Numerical results}
\subsection{Bounds}\label{experiments-bound}
We first construct bounds on the minimum distances of self-orthogonal codes over $\mathbb{F}_2$ with dimension $9 \leq k \leq 15$. Grassl's table~\cite{grassl} provides bounds for $d(n,k)$. Clearly, this gives an upper bound for $d_{SO}(n,k)$. Shi et al. gave a bound for self-orthogonal codes using residual codes.
\begin{lemma}\cite{bound}
    Let $\mathcal{C}$ be a self-orthogonal $[n,k,d]$ code over $\mathbb{F}_2$ and let $\mathbf{c}$ be a codeword of weight $w < 2d$, where $k \geq 2$. Then the residual code $\mathrm{Res}(\mathcal{C},\mathbf{c})$ of $\mathcal{C}$ is an even code with parameters
    \[
    \left[n-w, k-1, d' \geq 2 \left\lceil\frac{d}{2}-\frac{w}{4}\right\rceil\right].
    \] In particular, if $\mathbf{c}$ has weight $d$, where $d \equiv 2 \pmod4$, then $\mathrm{Res}(\mathcal{C},\mathbf{c})$ has parameters \[
    \left[ n-d, k-1, d' \geq \frac{d}{2}+1 \right].
    \]
\end{lemma}
Given $n$ and $k \geq 2$, let $d=d(n,k)$. Assume that there is a self-orthogonal $[n,k,d]$ code over $\mathbb{F}_2$. For a codeword $\mathbf{c} \in \mathcal{C}$ of weight $d$, we consider the residual code $\mathrm{Res}(\mathcal{C},\mathbf{c})$ with parameters \[
\left[n-d, k-1, d' \geq 2 \left\lceil\frac{d}{4}\right\rceil\right].
\]
If no even code with these residual parameters exists, then $d_{SO}(n,k)$ cannot be equal to $d(n,k)$. Thus, \[
d_{SO}(n,k) \leq d(n,k)-2.
\]
Recently, Li et al.~\cite{bound2} showed that self-orthogonal codes with certain parameters do not exist.
\begin{enumerate}
    \item There are no binary self-orthogonal $[2^k-3, k, 2^{k-1}-2]$ and $[2^{k-1}-2, k, 2^{k-2}-2]$ codes for any $k \geq 4$.
    \item There are no binary self-orthogonal $[2^k-2^u-2, k, 2^{k-1}-2^{u-1}-2]$ codes for any $k \geq u+2 \geq 5$.
    \item There are no binary self-orthogonal $[2^k-1, k+1, 2^{k-1}-2]$ codes for $k \geq 5$.
\end{enumerate}
Combining Grassl's table~\cite{grassl} with the residual-code criterion of Shi et al.~\cite{bound} and the nonexistence results of Li et al.~\cite{bound2}, we derive upper bounds on $d_{SO}(n,k)$ for binary self-orthogonal $[n,k]$ codes in the range $20\leq n\leq 256$ and $9\leq k\leq 15$, which is presented in Table~\ref{tab:full-so-upper-bounds}. We denote by $\overline{d}_{SO}(n,k)$ the resulting upper bound on $d_{SO}(n,k)$.
For $k=9$ and $k=10$, the values of $d_{SO}(n,k)$ are already known for $20\leq n\leq 40$ in~\cite{optimal-so-bound}. Therefore, in these cases, we start from $n=41$.
\subsection{Experiments}
In this section, we describe our general methodology. Given a generator matrix $G$ of a binary linear code $\mathcal{C}$, our goal is to construct a shortest self-orthogonal embedding of $\mathcal{C}$ whose minimum distance is the largest among the embeddings found by our methodology.
\begin{enumerate}
    \item We aim to obtain self-orthogonal codes with large minimum distance of parameters in the range
    \[
    20\leq n\leq 256
    \qquad\text{and}\qquad
    9\leq k\leq 15.
    \]
    We use base codes from the MAGMA BKLC (Best Known Linear Codes) library and apply shortest self-orthogonal embedding methods.
    \item Let $\mathcal{C}$ be an $[n,k,d]$ code over $\mathbb{F}_2$ with generator matrix $G$, obtained from the MAGMA BKLC database, and assume that $\mathcal{C}$ is not self-orthogonal. Let $\ell=\dim\Hull(\mathcal{C})$. By Theorem~\ref{akklw-2025}, the number $m$ of appended columns in a shortest self-orthogonal embedding is $k-\ell$ if $\mathcal{C}$ is an odd-like code and $k-\ell+1$ if $\mathcal{C}$ is an even code.
    \item In our genetic algorithm, we use $\mathcal{T} = \{T_\mathbf{u}~|~\wt(\mathbf{u})=4~for~\mathbf{u}\in\mathbb{F}_2^m\}$ as genes. However, by Theorem~\ref{orthogonal-group-generator}, the orthogonal group of order $m$ is generated by $\mathcal{T}$ only for $m \geq 5$. Since $O(m,2)$ is equal to the permutation group $\mathcal{P}_m$ of order $m$ if $m\leq3$ and $O(4,2)$ is not solely generated by $\mathcal{T}$, GA is not feasible when $m \leq 4$. Hence, we implement GA for $m \geq 5$.
    \item \textbf{Genetic Algorithm. }The initial feasible matrix $B_0$ is computed using random search. Once the initial matrix $B_0$ is found, the genetic algorithm is run for at most $120$ generations, and the initial random chromosome length is bounded above by $3m$.
    \begin{enumerate}
        \item Initialization : The initial population consists of $120$ randomly generated chromosomes. Within the upper bound of chromosome length at most $3m$, the length of a chromosome is randomly selected. Genes are selected from $\mathcal{T}$ for predetermined chromosome length.

        \item Selection : Parent selection is performed by tournament selection with tournament size $5$. The top $2$ chromosomes are carried over from the previous generation as elites and $59$ pairs of parents are selected by tournament selection. In the selection phase, we use the fitness function $F(\chi)=\left(d_{\chi},-A_{d_{\chi}}(\mathcal{C}_{\chi}),-|\chi|\right)$ to evaluate each chromosome.

        \item Crossover : The crossover probability is set to $0.8$. When crossover is applied, random crossover is chosen with probability $0.45$, and guided crossover is chosen with probability $0.55$. Regarding guided crossover, we select $4$ prefix chromosomes from the base chromosome and $3$ suffix chromosomes from the secondary chromosome and choose $2$ offspring chromosomes which satisfy the length constraint $|\chi|\leq 4m$ and have the largest fitness values.

        \item Mutation : The mutation probability is set to $0.45$. When mutation is applied, the relative weights of the mutation operations are given by $0.29$ for replace, $0.29$ for insert, $0.21$ for delete, $0.21$ for swap. Moreover, the maximum allowable chromosome length is bounded above by $4m$.

        \item Local search : each offspring is refined by a local search of depth $8$ (Subsection~\ref{so-embedding-algorithm}).
    \end{enumerate}
    Finally, if the best fitness value does not improve for $20$ consecutive generations, a restart is performed by replacing $50\%$ of the population with newly generated chromosomes.
\begin{algorithm}[t]
\caption{Genetic Algorithm for Shortest Self-Orthogonal Embedding}
\label{alg:ga-so}
\begin{algorithmic}[1]
\Require A BKLC code $\mathcal{C}$ requiring $m\ (\ge 5)$ appended columns
\Ensure A self-orthogonal embedding of $\mathcal{C}$ with large minimum distance
\State Split the generator matrix of $\mathcal{C}$ into its hull part and complementary part $A$
\State Find an initial feasible matrix $B_0$ satisfying $AA^\top+B_0B_0^\top=\mathcal{O}$ by random search
\State Set the gene pool $\mathcal{T}=\{\,T_\mathbf{u}: \wt(\mathbf{u})=4,\ \mathbf{u}\in\mathbb{F}_2^m\,\}$
    \Statex \Comment{each chromosome is a gene sequence and yields an embedding by applying its genes to $B_0$}
\State Initialize $120$ chromosomes with random genes from $\mathcal{T}$, each of length at most $3m$
\For{at most $120$ generations}
    \State Carry over the best $2$ chromosomes as elites
    \For{$59$ parent pairs selected by tournament selection (size $5$)}
        \State With probability $0.8$ apply crossover (guided $0.55$, random $0.45$); otherwise keep the parents
        \State With probability $0.45$ mutate each offspring (replace $0.29$, insert $0.29$, delete $0.21$, swap $0.21$)
        \State Apply local search of depth $D$ to each offspring (append a random weight-$4$ transvection, keep on strict improvement)
        \State Add the two offspring, respecting the length bound $4m$
    \EndFor
    \State Rank chromosomes by the fitness function $F(\chi)=(d_{\chi},-A_{d_{\chi}}(\mathcal{C}_{\chi}),-|\chi|)$ in lexicographical order
    \If{the best value has not improved for $20$ generations}
        \State Replace $50\%$ of the population with new random chromosomes
    \EndIf
\EndFor
\State \Return the best embedding found
\end{algorithmic}
\end{algorithm}
\end{enumerate}
\subsection{Results}
We conducted experiments on a total of $469$ base codes with $m\geq 5$ obtained from the MAGMA BKLC database, where $m$ denotes the number of columns appended by a shortest self-orthogonal embedding. For each base code, our genetic algorithm searches for an embedding that maximizes the minimum distance of the resulting self-orthogonal code. In Tables~\ref{tab:so_embedding_gap0} and~\ref{tab:so_embedding_gap2}, $n$ denotes the length of the code obtained after applying a shortest self-orthogonal embedding, so that $[n,k]$ are the parameters of the embedded code with $9\leq k\leq 15$, and $d_{SO}^{GA}(n,k)$ denotes the largest minimum distance obtained by our method for these parameters. The upper bound on $d_{SO}(n,k)$ used for comparison is the one derived in Subsection~\ref{experiments-bound}.
We define the \emph{gap} of $(n,k)$ by
\[
\mathrm{gap}(n,k)=\overline{d}_{SO}(n,k)-d_{SO}^{GA}(n,k),
\]
that is, the amount by which the largest minimum distance achieved by our method falls short of the bound. In particular, a gap of $0$ means that the constructed self-orthogonal code is optimal, i.e., it meets the bound.
Among the $469$ base codes, our method attains the bound for $73$ of them and falls short of it by $2$ for another $151$. After discarding codes that coincide in both their parameters $[n,k]$ and their automorphism group order $|\mathrm{Aut}|$ (which we cannot distinguish here), these yield $201$ distinct self-orthogonal codes:
\begin{itemize}
    \item $66$ \emph{new optimal} self-orthogonal codes ($\mathrm{gap}=0$): since $d_{SO}^{GA}(n,k)$ equals the upper bound $\overline{d}_{SO}(n,k)$, they are optimal, and are listed in Table~\ref{tab:so_embedding_gap0}.
    \item $135$ self-orthogonal codes with $\mathrm{gap}=2$: here $\overline{d}_{SO}(n,k)=d_{SO}^{GA}(n,k)+2$, so optimality is not yet settled; they attain the best minimum distance obtained so far and are listed in Table~\ref{tab:so_embedding_gap2}.
\end{itemize}
All codes reported in both tables are obtained by the genetic algorithm.

\begingroup
\footnotesize
\renewcommand{\arraystretch}{1.3}
\setlength{\tabcolsep}{4pt}

\begin{longtable}{c|c|c||c|c|c||c|c|c||c|c|c}
\caption{New optimal self-orthogonal codes attaining the bound
($\mathrm{gap}=0$, i.e.,
$d_{SO}^{GA}(n,k)=\overline{d}_{SO}(n,k)$) with
$9\leq k\leq 15$, obtained from shortest self-orthogonal embeddings
using a genetic algorithm (GA). Here $|\mathrm{Aut}|$ denotes the
order of the permutation automorphism group; codes coinciding in
both $[n,k]$ and $|\mathrm{Aut}|$ are listed once.}
\label{tab:so_embedding_gap0}\\

\hline
$[n,k]$ & $d_{SO}^{GA}$ & $|\mathrm{Aut}|$
& $[n,k]$ & $d_{SO}^{GA}$ & $|\mathrm{Aut}|$
& $[n,k]$ & $d_{SO}^{GA}$ & $|\mathrm{Aut}|$
& $[n,k]$ & $d_{SO}^{GA}$ & $|\mathrm{Aut}|$ \\
\hline
\endfirsthead

\multicolumn{12}{c}{}\\
\hline
$[n,k]$ & $d_{SO}^{GA}$ & $|\mathrm{Aut}|$
& $[n,k]$ & $d_{SO}^{GA}$ & $|\mathrm{Aut}|$
& $[n,k]$ & $d_{SO}^{GA}$ & $|\mathrm{Aut}|$
& $[n,k]$ & $d_{SO}^{GA}$ & $|\mathrm{Aut}|$ \\
\hline
\endhead

\multicolumn{12}{r}{}\\
\endfoot

\endlastfoot

[43,9] & 16 & 1
& [249,10] & 120 & 8
& [27,12] & 8 & 80640
& [30,14] & 8 & 30 \\
\hline
[49,9] & 20 & 1
& [22,11] & 6 & 887040
& [28,12] & 8 & 22464
& [31,14] & 8 & 96 \\
\hline
[50,9] & 20 & 1
& [23,11] & 8 & 10200960
& [29,12] & 8 & 67392
& [31,14] & 8 & 2 \\
\hline
[159,9] & 76 & 8
& [24,11] & 8 & 10200960
& [30,12] & 8 & 4
& [32,14] & 8 & 3584 \\
\hline
[160,9] & 76 & 1
& [25,11] & 8 & 20401920
& [44,12] & 16 & 32
& [73,14] & 28 & 1 \\
\hline
[208,9] & 100 & 1
& [27,11] & 8 & 144
& [46,12] & 16 & 32
& [81,14] & 32 & 128 \\
\hline
[215,9] & 104 & 2
& [28,11] & 8 & 384
& [47,12] & 16 & 32
& [127,14] & 56 & 889 \\
\hline
[216,9] & 104 & 1
& [188,11] & 88 & 4
& [125,12] & 56 & 1
& [31,15] & 8 & 2688 \\
\hline
[223,9] & 108 & 1
& [188,11] & 88 & 1
& [233,12] & 110 & 1
& [32,15] & 8 & 720 \\
\hline
[253,9] & 124 & 3600
& [189,11] & 88 & 2
& [26,13] & 6 & 11232
& [32,15] & 8 & 8 \\
\hline
[254,9] & 124 & 1440
& [189,11] & 88 & 1
& [30,13] & 8 & 8
& [33,15] & 8 & 930 \\
\hline
[42,10] & 16 & 64
& [191,11] & 90 & 2
& [31,13] & 8 & 4
& [81,15] & 32 & 7 \\
\hline
[49,10] & 20 & 2
& [251,11] & 120 & 1
& [39,13] & 12 & 1
& [82,15] & 32 & 64 \\
\hline
[50,10] & 20 & 3
& [252,11] & 120 & 1
& [78,13] & 32 & 128
& [82,15] & 32 & 1 \\
\hline
[58,10] & 24 & 42
& [24,12] & 8 & 244823040
& [79,13] & 32 & 128
& [128,15] & 56 & 113792 \\
\hline
[66,10] & 28 & 336
& [26,12] & 8 & 11232
& [126,13] & 56 & 14
& & & \\
\hline
[130,10] & 60 & 64
& [27,12] & 8 & 11232
& [127,13] & 56 & 14
& & & \\
\hline
\end{longtable}
\endgroup



\begingroup
\footnotesize
\renewcommand{\arraystretch}{1.3}
\setlength{\tabcolsep}{4pt}

\begin{longtable}{c|c|c||c|c|c||c|c|c||c|c|c}
\caption{Self-orthogonal codes lying two below the bound
($\mathrm{gap}=2$, i.e.,
$d_{SO}^{GA}(n,k)=\overline{d}_{SO}(n,k)-2$) with
$9\leq k\leq 15$, obtained from shortest self-orthogonal embeddings
using a genetic algorithm (GA). Here $|\mathrm{Aut}|$ denotes the
order of the permutation automorphism group; codes coinciding in
both $[n,k]$ and $|\mathrm{Aut}|$ are listed once.}
\label{tab:so_embedding_gap2}\\

\hline
$[n,k]$ & $d_{SO}^{GA}$ & $|\mathrm{Aut}|$
& $[n,k]$ & $d_{SO}^{GA}$ & $|\mathrm{Aut}|$
& $[n,k]$ & $d_{SO}^{GA}$ & $|\mathrm{Aut}|$
& $[n,k]$ & $d_{SO}^{GA}$ & $|\mathrm{Aut}|$ \\
\hline
\endfirsthead

\multicolumn{12}{c}{}\\
\hline
$[n,k]$ & $d_{SO}^{GA}$ & $|\mathrm{Aut}|$
& $[n,k]$ & $d_{SO}^{GA}$ & $|\mathrm{Aut}|$
& $[n,k]$ & $d_{SO}^{GA}$ & $|\mathrm{Aut}|$
& $[n,k]$ & $d_{SO}^{GA}$ & $|\mathrm{Aut}|$ \\
\hline
\endhead

\multicolumn{12}{r}{}\\
\endfoot

\endlastfoot

[44,9] & 16 & 2 & [187,11] & 86 & 1
& [38,13] & 10 & 1 & [82,14] & 32 & 1 \\
\hline
[52,9] & 20 & 1 & [190,11] & 88 & 1
& [40,13] & 12 & 1 & [83,14] & 32 & 2 \\
\hline
[83,9] & 36 & 1 & [253,11] & 120 & 2
& [41,13] & 12 & 1 & [147,14] & 64 & 1 \\
\hline
[84,9] & 36 & 1 & [25,12] & 6 & 120960
& [46,13] & 14 & 2 & [188,14] & 84 & 1 \\
\hline
[148,9] & 68 & 4 & [31,12] & 8 & 4
& [50,13] & 16 & 1 & [189,14] & 84 & 1 \\
\hline
[149,9] & 70 & 1 & [32,12] & 8 & 1
& [73,13] & 28 & 2 & [191,14] & 86 & 6 \\
\hline
[150,9] & 70 & 1 & [33,12] & 8 & 1
& [79,13] & 30 & 1 & [192,14] & 86 & 6 \\
\hline
[189,9] & 90 & 1 & [39,12] & 12 & 1
& [80,13] & 30 & 1 & [204,14] & 92 & 3 \\
\hline
[210,9] & 100 & 1 & [40,12] & 12 & 1
& [114,13] & 48 & 2 & [210,14] & 96 & 3 \\
\hline
[248,9] & 120 & 98304 & [41,12] & 12 & 2
& [146,13] & 64 & 20 & [211,14] & 96 & 3 \\
\hline
[250,9] & 120 & 9216 & [44,12] & 14 & 1
& [155,13] & 68 & 15 & [228,14] & 104 & 1 \\
\hline
[251,9] & 120 & 6 & [78,12] & 30 & 1
& [159,13] & 70 & 3 & [229,14] & 104 & 1 \\
\hline
[43,10] & 14 & 1 & [124,12] & 54 & 1
& [176,13] & 78 & 14 & [30,15] & 6 & 576 \\
\hline
[84,10] & 36 & 1 & [158,12] & 70 & 408
& [190,13] & 86 & 2 & [34,15] & 8 & 16128 \\
\hline
[85,10] & 36 & 1 & [188,12] & 86 & 1
& [191,13] & 86 & 6 & [34,15] & 8 & 1 \\
\hline
[149,10] & 68 & 2 & [189,12] & 86 & 1
& [203,13] & 92 & 3 & [35,15] & 8 & 1 \\
\hline
[167,10] & 78 & 1 & [201,12] & 92 & 1
& [210,13] & 96 & 2 & [36,15] & 8 & 2 \\
\hline
[188,10] & 88 & 2 & [206,12] & 94 & 24
& [224,13] & 102 & 8 & [37,15] & 8 & 6 \\
\hline
[189,10] & 88 & 2 & [207,12] & 96 & 1
& [227,13] & 104 & 2 & [39,15] & 10 & 1 \\
\hline
[190,10] & 90 & 2 & [229,12] & 106 & 1
& [28,14] & 6 & 11232 & [41,15] & 10 & 1 \\
\hline
[191,10] & 90 & 2 & [230,12] & 108 & 1
& [33,14] & 8 & 1152 & [48,15] & 14 & 1 \\
\hline
[244,10] & 116 & 8 & [231,12] & 108 & 1
& [34,14] & 8 & 1 & [50,15] & 14 & 1 \\
\hline
[248,10] & 118 & 8 & [232,12] & 108 & 1
& [35,14] & 8 & 2 & [53,15] & 16 & 32 \\
\hline
[251,10] & 120 & 32 & [251,12] & 118 & 3
& [35,14] & 8 & 3 & [58,15] & 18 & 1 \\
\hline
[31,11] & 8 & 1 & [252,12] & 118 & 3
& [38,14] & 10 & 1 & [74,15] & 26 & 1 \\
\hline
[37,11] & 12 & 1 & [27,13] & 6 & 138240
& [40,14] & 10 & 1 & [77,15] & 28 & 1 \\
\hline
[39,11] & 12 & 1 & [28,13] & 6 & 138240
& [41,14] & 12 & 1 & [83,15] & 32 & 1 \\
\hline
[41,11] & 14 & 1 & [28,13] & 6 & 1152
& [46,14] & 14 & 1 & [84,15] & 32 & 2 \\
\hline
[42,11] & 14 & 64 & [29,13] & 6 & 276480
& [47,14] & 14 & 1 & [85,15] & 32 & 1 \\
\hline
[43,11] & 14 & 32 & [30,13] & 6 & 829440
& [48,14] & 14 & 1 & [149,15] & 64 & 128 \\
\hline
[46,11] & 16 & 1 & [33,13] & 8 & 1
& [56,14] & 18 & 1 & [190,15] & 84 & 2 \\
\hline
[104,11] & 44 & 1 & [34,13] & 8 & 2
& [64,14] & 22 & 1 & [193,15] & 86 & 6 \\
\hline
[119,11] & 52 & 1 & [34,13] & 8 & 1
& [74,14] & 28 & 1 & [194,15] & 86 & 6 \\
\hline
[121,11] & 54 & 1 & [35,13] & 8 & 1
& [80,14] & 30 & 64 & & & \\
\hline
\end{longtable}
\endgroup
The generator matrices of resulting self-orthogonal embeddings are available at~\cite{github}.

\subsection{Running time}
To confirm the analysis of Subsection~\ref{sec:complexity} empirically, Table~\ref{tab:runtime} reports measured running times for one representative base code in each dimension $k=9,\ldots,15$. As predicted, the cost of a single fitness evaluation is governed by the number of cosets $2^{r}$: for the codes with $r=4$ it is only about $3$--$5\,\mu$s, whereas for the code with $r=10$ it rises to about $277\,\mu$s, i.e.\ roughly $60$--$80$ times larger, in agreement with the factor $2^{10-4}=64$. The one-time coset preprocessing and the total running time remain well under a second in all cases.

\begin{table}[htbp]
\centering
\caption{Measured running times of the genetic algorithm for one representative base code per dimension. Here $r=k-\dim\Hull(\mathcal{C})$, $m$ is the number of appended columns, and the average fitness-evaluation time confirms the $O(2^{r})$ cost of Subsection~\ref{sec:complexity}.}
\label{tab:runtime}
\small
\renewcommand{\arraystretch}{1.2}
\setlength{\tabcolsep}{5pt}
\begin{tabular}{c|c|c|r|r|r|r|r}
\Xhline{2\arrayrulewidth}
$[n,k]$ & $r$ & $m$ & Gen. & Evals & Avg.\ eval ($\mu$s) & Prep.\ (ms) & Total (s) \\
\Xhline{2\arrayrulewidth}
$[44,9]$   & $4$  & $5$  & $1$   & $120$    & $3.30$   & $0.045$ & $0.034$ \\
\hline
$[185,10]$ & $4$  & $5$  & $120$ & $128{,}362$ & $3.32$   & $0.086$ & $0.617$ \\
\hline
$[177,11]$ & $10$ & $11$ & $3$   & $3{,}255$   & $277.33$ & $0.568$ & $0.957$ \\
\hline
$[120,12]$ & $4$  & $5$  & $1$   & $120$    & $5.14$   & $0.292$ & $0.035$ \\
\hline
$[186,13]$ & $4$  & $5$  & $120$ & $128{,}185$ & $3.43$   & $0.625$ & $0.639$ \\
\hline
$[186,14]$ & $4$  & $5$  & $120$ & $128{,}094$ & $3.30$   & $1.179$ & $0.617$ \\
\hline
$[77,15]$  & $5$  & $5$  & $1$   & $120$    & $5.30$   & $2.001$ & $0.028$ \\
\Xhline{2\arrayrulewidth}
\end{tabular}
\end{table}
\subsection{Ablation study on guided crossover}
\label{sec:ablation-guided}

To assess the contribution of the guided crossover, we compare two configurations of the genetic algorithm on the $201$ codes that admit a self-orthogonal embedding attaining the bound (gap $0$, $66$ codes) or lying $2$ below it (gap $2$, $135$ codes). In the \emph{guided} configuration, crossover is applied with probability $0.8$; when applied, guided crossover is chosen with probability $0.55$ and random crossover with probability $0.45$. In the \emph{plain} configuration, crossover is likewise applied with probability $0.8$ but is always random ($100\%$ random crossover). All other operators and parameters are identical. To keep the running times of the two configurations comparable, the guided configuration is run for $30$ generations, whereas the plain configuration is run for $50$ generations.

Table~\ref{tab:ablation_guided} summarizes the outcome. The guided configuration attains all $66$ bound-meeting codes, whereas the plain configuration misses $2$ of them even with $20$ additional generations. On the gap-$2$ codes the two configurations are essentially tied ($132$ versus $133$ out of $135$), so guided crossover is not uniformly superior. Nevertheless, its ability to reach every bound-attaining target within a shorter generation budget indicates that it explores diverse regions of the search space and thereby escapes local optima to reach the optimal codes. For this reason we adopt guided crossover as a standard operator.

The running times of the two configurations are nearly identical, as reported in Table~\ref{tab:ablation_guided}: the per-case medians are close ($0.80$\,s versus $0.95$\,s) and the total wall-clock times differ by less than $4\%$ ($55$\,min\,$28$\,s versus $57$\,min\,$22$\,s).

\begin{table}[htbp]
\centering
\caption{Ablation of guided crossover on the $201$ codes ($66$ with gap $0$, $135$ with gap $2$). The columns ``gap $0$'' and ``gap $2$'' count the codes for which the configuration reaches the corresponding target, and the remaining columns report per-case running time.}
\label{tab:ablation_guided}
\small
\renewcommand{\arraystretch}{1.3}
\setlength{\tabcolsep}{5pt}
\begin{tabular}{l|c|c|r|r|r|r}
\Xhline{2\arrayrulewidth}
Configuration & gap $0$ ($/66$) & gap $2$ ($/135$) & Total time & Mean & Median & Max \\
\Xhline{2\arrayrulewidth}
Guided ($30$ gen.) & $66$ & $132$ & $3328.07$\,s ($\approx 55$\,m\,$28$\,s) & $16.56$\,s & $0.80$\,s & $249.28$\,s \\
\hline
Plain ($50$ gen.)  & $64$ & $133$ & $3442.18$\,s ($\approx 57$\,m\,$22$\,s) & $17.13$\,s & $0.95$\,s & $261.50$\,s \\
\Xhline{2\arrayrulewidth}
\end{tabular}
\end{table}

\subsection{Comparison with random search}
\label{sec:ga-vs-random}

To justify the use of a genetic algorithm, we compare it against random search under an equal time budget. Since the genetic algorithm already employs a local search of depth $8$, we consider two random baselines: random search with the same local search of depth $8$ (\emph{Random\,+\,LS(8)}) and pure random search with no local search (\emph{Pure random}). The genetic algorithm is the guided crossover version. For each code we run the genetic algorithm, record its wall-clock time, and grant each random baseline exactly that time as its budget. Because random search has no notion of generations, a time-based budget is the appropriate basis for a fair comparison. All three methods search the same feasible set of embeddings for the $201$ codes.

Table~\ref{tab:ga_vs_random} reports the results. Within the shared time budget, the genetic algorithm attains all $66$ bound-meeting (gap $0$) codes, whereas Random\,+\,LS(8) attains only $57$ and Pure random only $49$. The genetic algorithm also finds the largest number of gap-$\leq 2$ codes ($198$ out of $201$) and the smallest mean gap. Moreover, it is never outperformed on any single case: it strictly improves on Random\,+\,LS(8) in $18$ codes and on Pure random in $28$ codes, with no losses.

The comparison also clarifies the role of local search: adding it to random search raises the number of optimal codes from $49$ to $57$, yet this is still well below the $66$ obtained by the genetic algorithm. Hence local refinement alone is not sufficient, and the evolutionary operators contribute beyond it. This is further supported by the number of candidates evaluated within the budget: Random\,+\,LS(8) evaluates about $6{,}200$ candidates per case (median) and Pure random about $39{,}000$; despite exploring an order of magnitude more candidates, pure random search finds the fewest optimal codes. Thus raw sampling volume does not compensate for the lack of guided evolution, which justifies the use of the genetic algorithm.


\begin{table}[htbp]
\centering
\caption{Genetic algorithm versus random search under an equal per-case time budget (the genetic algorithm running time), over the $201$ codes. ``gap $0$'' counts the bound-attaining (optimal) codes found, and ``gap $\leq 2$'' counts the codes within $2$ of the bound.}
\label{tab:ga_vs_random}
\small
\renewcommand{\arraystretch}{1.3}
\setlength{\tabcolsep}{6pt}
\begin{tabular}{l|c|c|c}
\Xhline{2\arrayrulewidth}
Method & gap $0$ ($/66$) & gap $\leq 2$ ($/201$) & Mean gap \\
\Xhline{2\arrayrulewidth}
Genetic algorithm (guided crossover, LS depth $8$) & $66$ & $198$ & $1.37$ \\
\hline
Random search $+$ LS depth $8$ & $57$ & $189$ & $1.55$ \\
\hline
Pure random search & $49$ & $187$ & $1.65$ \\
\Xhline{2\arrayrulewidth}
\end{tabular}
\end{table}
\FloatBarrier


\section{Conclusion}
In this paper, we developed a constraint-preserving genetic algorithm for constructing self-orthogonal codes through shortest self-orthogonal embeddings. On the theoretical side, we reformulated the search for a shortest self-orthogonal embedding of an $[n,k]$ code as the problem of finding a block $B$ satisfying the feasibility condition $BB^{T}=AA^{T}$, and we used the fact that all shortest embeddings are obtained by the action of orthogonal group $O(m,2)$. Building on the result that $O(m,2)$ is generated by the weight-four transvections for $m\geq5$, we encoded each chromosome as a sequence of such transvections; in this representation every chromosome yields a feasible embedding, so the crossover, mutation, and local-search operators automatically preserve the self-orthogonality constraint. 
Using this framework, we found $66$ new optimal self-orthogonal codes and $135$ further self-orthogonal codes of best known minimum distance in dimensions $9\leq k\leq15$, whereas previous computational searches were limited to dimensions up to $8$. An ablation study and a comparison with random search under an equal time budget further confirm that the genetic algorithm, and in particular the guided crossover, is an efficient and effective method for this problem.

\section*{Acknowledgement}
This research (J.-L. Kim) was supported in part by the BK21 FOUR (Fostering Outstanding Universities for Research) funded by the Ministry of Education (MOE, Korea), National Research Foundation of Korea (NRF) under Grant No. 4120240415042, Basic Science Research Program through the National Research Foundation of Korea (NRF) funded by the Ministry of Science and ICT under Grant No. RS-2025-24534992 and Global - Learning \& Academic research institution for Master’s·PhD students, and Postdocs(LAMP) Program of the National Research Foundation of Korea(NRF) grant funded by the Ministry of Education(No. RS-2024-00441954).

\section*{Appendix}
Here, we present the upper bounds on the minimum distances of self-orthogonal codes over $\mathbb{F}_2$ with dimensions $9\le k\le 15$ and lengths $20\le n\le 256$. In Table~\ref{tab:full-so-upper-bounds}, each entry represents an upper bound on the minimum distance of the corresponding self-orthogonal $[n,k]$ linear code over $\mathbb{F}_2$.


\begingroup
\sffamily\footnotesize
\setlength{\tabcolsep}{1pt}
\renewcommand{\arraystretch}{0.98}

\refstepcounter{table}
\label{tab:full-so-upper-bounds}
\par\addvspace{6pt}
\noindent\textbf{\tablename~\thetable}\par
\noindent Upper bounds on the minimum distances of binary self-orthogonal codes with $9\leq k\leq15$ and $20\leq n\leq256$.\par\vspace{4pt}
\addtocounter{table}{-1}

\begin{longtable}{>{\centering\arraybackslash}p{0.68cm}|*{7}{>{\centering\arraybackslash}p{0.76cm}}@{\hspace{0.7em}\vrule width 0.3pt\hspace{0.7em}}>{\centering\arraybackslash}p{0.68cm}|*{7}{>{\centering\arraybackslash}p{0.76cm}}}
\toprule
\multicolumn{1}{c|}{$n\backslash k$} & $9$ & $10$ & $11$ & $12$ & $13$ & $14$ & $15$ & \multicolumn{1}{c|}{$n\backslash k$} & $9$ & $10$ & $11$ & $12$ & $13$ & $14$ & $15$ \\
\midrule
\endfirsthead
\multicolumn{16}{c}{\tablename~\thetable{} -- continued from previous page}\\
\toprule
\multicolumn{1}{c|}{$n\backslash k$} & $9$ & $10$ & $11$ & $12$ & $13$ & $14$ & $15$ & \multicolumn{1}{c|}{$n\backslash k$} & $9$ & $10$ & $11$ & $12$ & $13$ & $14$ & $15$ \\
\midrule
\endhead
\midrule
\multicolumn{16}{r}{Continued on next page}\\
\endfoot
\bottomrule
\endlastfoot
20 &  &  & 4 & 4 & 4 & 4 & 2 & 139 & 64 & 64 & 64 & 64 & 64 & 62 & 62 \\
21 &  &  & 6 & 4 & 4 & 4 & 4 & 140 & 66 & 64 & 64 & 64 & 64 & 62 & 62 \\
22 &  &  & 6 & 6 & 4 & 4 & 4 & 141 & 66 & 66 & 64 & 64 & 64 & 64 & 62 \\
23 &  &  & 8 & 6 & 4 & 4 & 4 & 142 & 68 & 66 & 66 & 64 & 64 & 64 & 64 \\
24 &  &  & 8 & 8 & 6 & 4 & 4 & 143 & 68 & 68 & 66 & 64 & 64 & 64 & 64 \\
25 &  &  & 8 & 8 & 6 & 6 & 4 & 144 & 68 & 68 & 68 & 66 & 64 & 64 & 64 \\
26 &  &  & 8 & 8 & 6 & 6 & 6 & 145 & 68 & 68 & 68 & 66 & 66 & 64 & 64 \\
27 &  &  & 8 & 8 & 8 & 6 & 6 & 146 & 68 & 68 & 68 & 68 & 66 & 66 & 64 \\
28 &  &  & 8 & 8 & 8 & 8 & 6 & 147 & 70 & 68 & 68 & 68 & 68 & 66 & 66 \\
29 &  &  & 8 & 8 & 8 & 8 & 6 & 148 & 70 & 70 & 68 & 68 & 68 & 66 & 66 \\
30 &  &  & 10 & 8 & 8 & 8 & 8 & 149 & 72 & 70 & 70 & 68 & 68 & 68 & 66 \\
31 &  &  & 10 & 10 & 8 & 8 & 8 & 150 & 72 & 72 & 70 & 68 & 68 & 68 & 68 \\
32 &  &  & 12 & 10 & 10 & 8 & 8 & 151 & 72 & 72 & 70 & 70 & 68 & 68 & 68 \\
33 &  &  & 12 & 10 & 10 & 10 & 8 & 152 & 72 & 72 & 72 & 70 & 70 & 68 & 68 \\
34 &  &  & 12 & 12 & 10 & 10 & 10 & 153 & 72 & 72 & 72 & 70 & 70 & 70 & 68 \\
35 &  &  & 12 & 12 & 10 & 10 & 10 & 154 & 72 & 72 & 72 & 72 & 70 & 70 & 70 \\
36 &  &  & 12 & 12 & 12 & 10 & 10 & 155 & 74 & 72 & 72 & 72 & 70 & 70 & 70 \\
37 &  &  & 14 & 12 & 12 & 12 & 10 & 156 & 74 & 74 & 72 & 72 & 72 & 70 & 70 \\
38 &  &  & 14 & 14 & 12 & 12 & 12 & 157 & 76 & 74 & 72 & 72 & 72 & 72 & 70 \\
39 &  &  & 14 & 14 & 12 & 12 & 12 & 158 & 76 & 76 & 74 & 72 & 72 & 72 & 72 \\
40 & 16 & 16 & 14 & 14 & 14 & 12 & 12 & 159 & 76 & 76 & 74 & 74 & 72 & 72 & 72 \\
41 & 16 & 16 & 16 & 14 & 14 & 14 & 12 & 160 & 76 & 76 & 76 & 74 & 74 & 72 & 72 \\
42 & 16 & 16 & 16 & 16 & 14 & 14 & 14 & 161 & 76 & 76 & 76 & 74 & 74 & 72 & 72 \\
43 & 16 & 16 & 16 & 16 & 14 & 14 & 14 & 162 & 78 & 76 & 76 & 76 & 74 & 74 & 72 \\
44 & 18 & 16 & 16 & 16 & 16 & 14 & 14 & 163 & 78 & 78 & 76 & 76 & 74 & 74 & 74 \\
45 & 18 & 18 & 16 & 16 & 16 & 16 & 14 & 164 & 80 & 78 & 76 & 76 & 76 & 74 & 74 \\
46 & 20 & 18 & 18 & 16 & 16 & 16 & 16 & 165 & 80 & 78 & 78 & 76 & 76 & 76 & 74 \\
47 & 20 & 20 & 18 & 16 & 16 & 16 & 16 & 166 & 80 & 80 & 78 & 76 & 76 & 76 & 76 \\
48 & 20 & 20 & 18 & 18 & 16 & 16 & 16 & 167 & 80 & 80 & 78 & 78 & 76 & 76 & 76 \\
49 & 20 & 20 & 20 & 18 & 18 & 16 & 16 & 168 & 80 & 80 & 80 & 78 & 78 & 76 & 76 \\
50 & 20 & 20 & 20 & 20 & 18 & 18 & 16 & 169 & 80 & 80 & 80 & 78 & 78 & 76 & 76 \\
51 & 22 & 20 & 20 & 20 & 20 & 18 & 18 & 170 & 80 & 80 & 80 & 80 & 78 & 78 & 76 \\
52 & 22 & 22 & 20 & 20 & 20 & 18 & 18 & 171 & 82 & 80 & 80 & 80 & 78 & 78 & 78 \\
53 & 22 & 22 & 22 & 20 & 20 & 20 & 18 & 172 & 82 & 82 & 80 & 80 & 80 & 78 & 78 \\
54 & 24 & 22 & 22 & 22 & 20 & 20 & 20 & 173 & 84 & 82 & 80 & 80 & 80 & 80 & 78 \\
55 & 24 & 24 & 22 & 22 & 20 & 20 & 20 & 174 & 84 & 84 & 82 & 80 & 80 & 80 & 80 \\
56 & 24 & 24 & 24 & 22 & 22 & 20 & 20 & 175 & 84 & 84 & 82 & 82 & 80 & 80 & 80 \\
57 & 24 & 24 & 24 & 22 & 22 & 20 & 20 & 176 & 84 & 84 & 84 & 82 & 80 & 80 & 80 \\
58 & 24 & 24 & 24 & 24 & 22 & 22 & 20 & 177 & 84 & 84 & 84 & 82 & 82 & 80 & 80 \\
59 & 26 & 24 & 24 & 24 & 24 & 22 & 22 & 178 & 86 & 84 & 84 & 84 & 82 & 82 & 80 \\
60 & 26 & 24 & 24 & 24 & 24 & 22 & 22 & 179 & 86 & 86 & 84 & 84 & 82 & 82 & 82 \\
61 & 26 & 26 & 24 & 24 & 24 & 24 & 22 & 180 & 88 & 86 & 84 & 84 & 84 & 82 & 82 \\
62 & 28 & 26 & 26 & 24 & 24 & 24 & 24 & 181 & 88 & 86 & 86 & 84 & 84 & 84 & 82 \\
63 & 28 & 28 & 26 & 26 & 24 & 24 & 24 & 182 & 88 & 88 & 86 & 84 & 84 & 84 & 84 \\
64 & 28 & 28 & 26 & 26 & 26 & 24 & 24 & 183 & 88 & 88 & 86 & 86 & 84 & 84 & 84 \\
65 & 28 & 28 & 28 & 26 & 26 & 24 & 24 & 184 & 88 & 88 & 88 & 86 & 84 & 84 & 84 \\
66 & 30 & 28 & 28 & 28 & 26 & 26 & 24 & 185 & 88 & 88 & 88 & 86 & 86 & 84 & 84 \\
67 & 30 & 28 & 28 & 28 & 28 & 26 & 26 & 186 & 90 & 88 & 88 & 88 & 86 & 86 & 84 \\
68 & 30 & 30 & 28 & 28 & 28 & 28 & 26 & 187 & 90 & 88 & 88 & 88 & 86 & 86 & 86 \\
69 & 30 & 30 & 28 & 28 & 28 & 28 & 26 & 188 & 90 & 90 & 88 & 88 & 88 & 86 & 86 \\
70 & 32 & 30 & 30 & 28 & 28 & 28 & 28 & 189 & 92 & 90 & 88 & 88 & 88 & 86 & 86 \\
71 & 32 & 32 & 30 & 30 & 28 & 28 & 28 & 190 & 92 & 92 & 90 & 88 & 88 & 88 & 86 \\
72 & 32 & 32 & 32 & 30 & 30 & 28 & 28 & 191 & 92 & 92 & 90 & 90 & 88 & 88 & 88 \\
73 & 32 & 32 & 32 & 30 & 30 & 28 & 28 & 192 & 92 & 92 & 92 & 90 & 88 & 88 & 88 \\
74 & 32 & 32 & 32 & 32 & 30 & 30 & 28 & 193 & 94 & 92 & 92 & 90 & 90 & 88 & 88 \\
75 & 32 & 32 & 32 & 32 & 32 & 30 & 30 & 194 & 94 & 92 & 92 & 92 & 90 & 90 & 88 \\
76 & 34 & 32 & 32 & 32 & 32 & 32 & 30 & 195 & 94 & 94 & 92 & 92 & 90 & 90 & 90 \\
77 & 34 & 34 & 32 & 32 & 32 & 32 & 30 & 196 & 96 & 94 & 92 & 92 & 90 & 90 & 90 \\
78 & 36 & 34 & 34 & 32 & 32 & 32 & 32 & 197 & 96 & 94 & 94 & 92 & 92 & 90 & 90 \\
79 & 36 & 36 & 34 & 34 & 32 & 32 & 32 & 198 & 96 & 96 & 94 & 92 & 92 & 92 & 90 \\
80 & 36 & 36 & 34 & 34 & 32 & 32 & 32 & 199 & 96 & 96 & 96 & 94 & 92 & 92 & 92 \\
81 & 36 & 36 & 36 & 34 & 34 & 32 & 32 & 200 & 96 & 96 & 96 & 94 & 92 & 92 & 92 \\
82 & 36 & 36 & 36 & 36 & 34 & 34 & 32 & 201 & 96 & 96 & 96 & 94 & 94 & 92 & 92 \\
83 & 38 & 36 & 36 & 36 & 36 & 34 & 34 & 202 & 96 & 96 & 96 & 96 & 94 & 94 & 92 \\
84 & 38 & 38 & 36 & 36 & 36 & 36 & 34 & 203 & 98 & 96 & 96 & 96 & 94 & 94 & 94 \\
85 & 40 & 38 & 38 & 36 & 36 & 36 & 34 & 204 & 98 & 96 & 96 & 96 & 96 & 94 & 94 \\
86 & 40 & 38 & 38 & 36 & 36 & 36 & 36 & 205 & 100 & 98 & 96 & 96 & 96 & 94 & 94 \\
87 & 40 & 40 & 38 & 38 & 36 & 36 & 36 & 206 & 100 & 98 & 98 & 96 & 96 & 96 & 94 \\
88 & 40 & 40 & 38 & 38 & 38 & 36 & 36 & 207 & 100 & 100 & 98 & 98 & 96 & 96 & 96 \\
89 & 40 & 40 & 40 & 38 & 38 & 38 & 36 & 208 & 100 & 100 & 100 & 98 & 96 & 96 & 96 \\
90 & 42 & 40 & 40 & 40 & 38 & 38 & 36 & 209 & 100 & 100 & 100 & 100 & 98 & 96 & 96 \\
91 & 42 & 40 & 40 & 40 & 40 & 38 & 38 & 210 & 102 & 100 & 100 & 100 & 98 & 98 & 96 \\
92 & 42 & 40 & 40 & 40 & 40 & 40 & 38 & 211 & 102 & 100 & 100 & 100 & 98 & 98 & 98 \\
93 & 44 & 42 & 40 & 40 & 40 & 40 & 38 & 212 & 104 & 102 & 100 & 100 & 100 & 98 & 98 \\
94 & 44 & 42 & 42 & 40 & 40 & 40 & 40 & 213 & 104 & 102 & 102 & 100 & 100 & 100 & 98 \\
95 & 44 & 44 & 42 & 42 & 40 & 40 & 40 & 214 & 104 & 104 & 102 & 100 & 100 & 100 & 100 \\
96 & 44 & 44 & 42 & 42 & 42 & 40 & 40 & 215 & 104 & 104 & 104 & 102 & 100 & 100 & 100 \\
97 & 44 & 44 & 44 & 42 & 42 & 40 & 40 & 216 & 104 & 104 & 104 & 102 & 102 & 100 & 100 \\
98 & 46 & 44 & 44 & 44 & 42 & 42 & 40 & 217 & 104 & 104 & 104 & 102 & 102 & 100 & 100 \\
99 & 46 & 44 & 44 & 44 & 44 & 42 & 42 & 218 & 106 & 104 & 104 & 104 & 102 & 102 & 100 \\
100 & 46 & 46 & 44 & 44 & 44 & 42 & 42 & 219 & 106 & 104 & 104 & 104 & 102 & 102 & 102 \\
101 & 48 & 46 & 46 & 44 & 44 & 44 & 42 & 220 & 108 & 106 & 104 & 104 & 102 & 102 & 102 \\
102 & 48 & 46 & 46 & 46 & 44 & 44 & 44 & 221 & 108 & 106 & 106 & 104 & 104 & 102 & 102 \\
103 & 48 & 48 & 46 & 46 & 44 & 44 & 44 & 222 & 108 & 108 & 106 & 104 & 104 & 104 & 102 \\
104 & 48 & 48 & 46 & 46 & 46 & 44 & 44 & 223 & 108 & 108 & 108 & 106 & 104 & 104 & 104 \\
105 & 48 & 48 & 48 & 46 & 46 & 46 & 44 & 224 & 108 & 108 & 108 & 106 & 104 & 104 & 104 \\
106 & 48 & 48 & 48 & 48 & 46 & 46 & 44 & 225 & 110 & 108 & 108 & 106 & 106 & 104 & 104 \\
107 & 48 & 48 & 48 & 48 & 48 & 46 & 46 & 226 & 110 & 108 & 108 & 108 & 106 & 106 & 104 \\
108 & 50 & 48 & 48 & 48 & 48 & 46 & 46 & 227 & 110 & 110 & 108 & 108 & 106 & 106 & 106 \\
109 & 50 & 50 & 48 & 48 & 48 & 48 & 46 & 228 & 112 & 110 & 110 & 108 & 106 & 106 & 106 \\
110 & 52 & 50 & 50 & 48 & 48 & 48 & 48 & 229 & 112 & 112 & 110 & 108 & 108 & 106 & 106 \\
111 & 52 & 52 & 50 & 48 & 48 & 48 & 48 & 230 & 112 & 112 & 110 & 110 & 108 & 108 & 106 \\
112 & 52 & 52 & 50 & 50 & 48 & 48 & 48 & 231 & 112 & 112 & 112 & 110 & 108 & 108 & 108 \\
113 & 52 & 52 & 52 & 50 & 50 & 48 & 48 & 232 & 112 & 112 & 112 & 110 & 108 & 108 & 108 \\
114 & 52 & 52 & 52 & 52 & 50 & 50 & 48 & 233 & 112 & 112 & 112 & 110 & 110 & 108 & 108 \\
115 & 54 & 52 & 52 & 52 & 52 & 50 & 50 & 234 & 114 & 112 & 112 & 112 & 110 & 110 & 108 \\
116 & 54 & 54 & 52 & 52 & 52 & 50 & 50 & 235 & 114 & 112 & 112 & 112 & 112 & 110 & 110 \\
117 & 56 & 54 & 54 & 52 & 52 & 52 & 50 & 236 & 116 & 114 & 112 & 112 & 112 & 112 & 110 \\
118 & 56 & 56 & 54 & 52 & 52 & 52 & 52 & 237 & 116 & 114 & 114 & 112 & 112 & 112 & 112 \\
119 & 56 & 56 & 54 & 54 & 52 & 52 & 52 & 238 & 116 & 116 & 114 & 112 & 112 & 112 & 112 \\
120 & 56 & 56 & 56 & 54 & 54 & 52 & 52 & 239 & 116 & 116 & 116 & 114 & 112 & 112 & 112 \\
121 & 56 & 56 & 56 & 54 & 54 & 54 & 52 & 240 & 116 & 116 & 116 & 114 & 112 & 112 & 112 \\
122 & 56 & 56 & 56 & 56 & 54 & 54 & 54 & 241 & 116 & 116 & 116 & 116 & 114 & 112 & 112 \\
123 & 58 & 56 & 56 & 56 & 56 & 54 & 54 & 242 & 118 & 116 & 116 & 116 & 114 & 114 & 112 \\
124 & 58 & 58 & 56 & 56 & 56 & 54 & 54 & 243 & 118 & 118 & 116 & 116 & 116 & 114 & 114 \\
125 & 60 & 58 & 56 & 56 & 56 & 56 & 54 & 244 & 120 & 118 & 116 & 116 & 116 & 114 & 114 \\
126 & 60 & 58 & 58 & 56 & 56 & 56 & 56 & 245 & 120 & 120 & 118 & 116 & 116 & 114 & 114 \\
127 & 60 & 60 & 58 & 58 & 56 & 56 & 56 & 246 & 120 & 120 & 118 & 118 & 116 & 116 & 114 \\
128 & 60 & 60 & 60 & 58 & 58 & 56 & 56 & 247 & 120 & 120 & 120 & 118 & 116 & 116 & 116 \\
129 & 60 & 60 & 60 & 58 & 58 & 56 & 56 & 248 & 122 & 120 & 120 & 120 & 118 & 116 & 116 \\
130 & 62 & 60 & 60 & 60 & 58 & 58 & 56 & 249 & 122 & 120 & 120 & 120 & 118 & 118 & 116 \\
131 & 62 & 60 & 60 & 60 & 60 & 58 & 58 & 250 & 122 & 120 & 120 & 120 & 120 & 118 & 116 \\
132 & 62 & 62 & 60 & 60 & 60 & 58 & 58 & 251 & 122 & 122 & 120 & 120 & 120 & 120 & 118 \\
133 & 64 & 62 & 62 & 60 & 60 & 60 & 58 & 252 & 124 & 122 & 120 & 120 & 120 & 120 & 118 \\
134 & 64 & 64 & 62 & 60 & 60 & 60 & 60 & 253 & 124 & 124 & 122 & 120 & 120 & 120 & 120 \\
135 & 64 & 64 & 62 & 62 & 60 & 60 & 60 & 254 & 124 & 124 & 122 & 120 & 120 & 120 & 120 \\
136 & 64 & 64 & 64 & 62 & 62 & 60 & 60 & 255 & 124 & 124 & 124 & 122 & 120 & 120 & 120 \\
137 & 64 & 64 & 64 & 62 & 62 & 62 & 60 & 256 & 128 & 124 & 124 & 122 & 120 & 120 & 120 \\
138 & 64 & 64 & 64 & 64 & 62 & 62 & 62 &  &  &  &  &  &  &  &  \\
\end{longtable}
\endgroup




\begin{thebibliography}{1}
\bibitem{akklw-2025}
J. An, N. Kaplan, J.-L. Kim, J. Luo and G. Wang,
\emph{Shortest self-orthogonal embeddings of binary linear codes},
arXiv:2511.05440, 2025.
\bibitem{akl-2026}
J. An, J.-L. Kim and S. Ling,
\emph{Embedding linear codes over $\mathbb{Z}_4$ into self-orthogonal codes}, 2026.
\bibitem{bn-2001}
T. Blackmore and G. H. Norton,
\emph{Matrix-product codes over $F_q$},
Appl. Algebra Engrg. Comm. Comput.,
12 : 477--500, 2001.
\bibitem{optimal-so-bound}
I. Bouyukliev, S. Bouyuklieva, T. A. Gulliver and P. R. J. {\"O}sterg{\aa}rd,
\emph{Classification of optimal binary self-orthogonal codes},
J. Combin. Math. Combin. Comput.,
59 : 33--87, 2006.
\bibitem{crss-1998}
A. R. Calderbank, E. M. Rains, P. W. Shor and N. J. A. Sloane,
\emph{Quantum error correction via codes over {GF}(4)},
IEEE Trans. Inform. Theory,
44(4) : 1369--1387, 1998.
\bibitem{d-1975}
P. Delsarte,
\emph{On subfield subcodes of modified Reed–Solomon codes},
IEEE Trans. Inform. Theory,
21(5) : 575--576, 1975.
\bibitem{d-2009}
Y. Ding,
\emph{Asymptotic bound on binary self-orthogonal codes},
Sci. China Ser. A,
52 : 631--638, 2009.
\bibitem{f-1966}
G. D. Forney Jr.,
\emph{Concatenated Codes},
MIT Press, 1965.
\bibitem{goldberg-1989}
D. E. Goldberg,
Genetic Algorithms in Search, Optimization, and Machine Learning,
Addison-Wesley, Reading, MA, 1989.
\bibitem{grassl}
M. Grassl,
Code Tables: Bounds on the parameters of various types of codes,
Available from: \url{https://www.codetables.de/},
Acessed on June, 2026.
\bibitem{h-2021}
M. Harada,
\emph{Construction of binary LCD codes, ternary LCD codes and quaternary Hermitian LCD codes},
Des. Codes Cryptogr.,
89 : 2299--2312, 2021.
\bibitem{h-1975}
J. H. Holland,
Adaptation in Natural and Artificial Systems: An Introductory Analysis with Applications to Biology, Control, and Artificial Intelligence,
MIT press, 1992.
\bibitem{Parametrization}
G. J. Janusz,
\emph{Parametrization of self-dual codes by orthogonal matrices},
Finite Fields Appl.,
13(3) : 450–-491, 2007.
\bibitem{k-2023}
J.-L. Kim,
\emph{Binary optimal linear codes with various hull dimensions and entanglement-assisted QECCs},
Comput. Appl. Math.
42 : 114, 2023.
\bibitem{kc-2022}
J.-L. Kim and W.-H. Choi,
\emph{Self-orthogonality matrix and Reed-Muller codes},
IEEE Trans. Inform. Theory,
68(11) : 7159--7164, 2022.
\bibitem{kkl-2021}
J.-L. Kim, Y.-H. Kim and N. Lee,
\emph{Embedding linear codes into self-orthogonal codes and their optimal minimum distances},
IEEE Trans. Inform. Theory,
67(6) : 3701--3707, 2021.
\bibitem{kl-2004}
J.-L. Kim and Y. Lee,
\emph{Euclidean and Hermitian self-dual {MDS} codes over large finite fields},
J. Combin. Theory Ser. A,
105(1) : 79--95, 2004.
\bibitem{kll-2026}
J.-L. Kim, H. Liu and J. Luo,
\emph{How to expand a self-orthogonal code},
Adv. Math. Commun.,
22 : 233--239, 2026.
\bibitem{lk-2015}
Y. Lee and J.-L. Kim,
\emph{An efficient construction of self-dual codes},
Bull. Korean Math. Soc.,
52(3), 915--923, 2015.
\bibitem{bound2}
S. Li, M. Shi, T. Helleseth and S. Ling,
\emph{Nonexistence of Several Infinite Families of Binary Self-Orthogonal Codes},
IIEEE Trans. Inform. Theory,
71(1) : 285--297, 2026.
\bibitem{github}
\leavevmode\raggedright
H. Lim,
Self-Orthogonal-LCD-Embedding-using-Genetic-Algorithm-and-Greedy-Search,
Available from: \url{https://github.com/haeuncreate/Self-Orthogonal-LCD-Embedding-using-Genetic-Algorithm-and-Greedy-Search.git},
2026
\bibitem{m-1969}
F. J. MacWilliams,
\emph{Orthogonal matrices over finite fields},
Amer. Math. Monthly,
76(2) : 152--164, 1969.
\bibitem{Introduction}
V. Pless,
Introduction to the Theory of Error-Correcting Codes,
Wiley-Interscience, New York, 1998.
\bibitem{ps-1975}
V. Pless and N. J. A. Sloane,
\emph{On the Classification and Enumeration of Self-Dual Codes},
J. Combin. Theory Ser. A,
18(3) : 313--335, 1975.
\bibitem{p-1960}
M. Plotkin,
\emph{Binary codes with specified minimum distance},
IRE Trans. Information Theory,
6(4) : 445--450, 1960.
\bibitem{mat-theory-comm}
C. E. Shannon,
\emph{A Mathematical Theory of Communication},
Bell System Tech. J.,
27(3) : 379--423, 1948.
\bibitem{bound}
M. Shi, S. Li, T. Helleseth and J.-L. Kim,
\emph{Binary self-orthogonal codes which meet the Griesmer bound or have optimal minimum distances},
J. Combin. Theory, Ser. A,
214 : 106027, 2025.
\bibitem{st-2026}
M. Shi and S. Tao,
\emph{Expanding self-orthogonal codes over a ring $\mathbb{Z}_4$ to self-dual codes and unimodular lattices},
Cryptogr. Commun., 2026.
\bibitem{s-1972}
N. J. A. Sloane, S. M. Reddy and C.-L. Chen,
\emph{New binary codes},
IEEE Trans. Inform. Theory,
18(4) : 503--510, 1972.
\bibitem{wl-2026}
J. Wang and J. Luo,
\emph{Shortest embeddings of linear codes with arbitrary hull dimension},
arXiv:2604.08843, 2026.
\bibitem{w-2026}
B. Won, J.-L. Kim and Y. Lee,
\emph{Modified genetic algorithms for constructing new quaternary hermitian LCD codes},
Applied Soft Computing, 2026.
\end{thebibliography}


\end{document}

